\documentclass[10pt,a4paper]{article}
\usepackage{amsmath,amssymb,amsthm,graphicx}
\usepackage{amsmath,amsfonts,amssymb,amsthm,graphicx}
\usepackage{hyperref}
\usepackage{thmtools,thm-restate}
\usepackage[capitalize,nameinlink]{cleveref}
\usepackage{enumerate}
\allowdisplaybreaks
\setcounter{MaxMatrixCols}{30}
\newtheorem{theorem}{Theorem}
\newtheorem{claim}[theorem]{Claim}
\newtheorem{corollary}[theorem]{Corollary}
\newtheorem{definition}[theorem]{Definition}
\newtheorem{lemma}[theorem]{Lemma}
\newtheorem{problem}[theorem]{Problem}
\newtheorem{proposition}[theorem]{Proposition}
\newcommand{\<}{\langle}
\renewcommand{\>}{\rangle}
\DeclareMathOperator{\E}{\mathbb E}
\newcommand{\cl}[1]{\mathsf{#1}}
\newcommand{\ApxCount}{\mathsf{ApxCount}_{N,w}}
\newcommand{\AndApxCount}{\mathsf{AND}_2\circ\mathsf{ApxCount}_{N,w}}
\newcommand{\para}[1]{\paragraph{#1}}
\newcommand{\id}{\mathbf 1}
\newcommand{\R}{\mathcal R}
\newcommand{\Reals}{\mathbb R}
\newcommand{\eps}{\varepsilon}
\title{Optimal approximate-counting bounds with trusted quantum samples, reflections and a QSampling unitary}
\author{Scott Aaronson, Robin Kothari, William Kretschmer, Justin Thaler}
\date{doi:10.4230/LIPIcs.CCC.2020.7; CC BY 3.0 Unported}
\begin{document}
\maketitle
\noindent\textit{Editorial note.} Editable transcription of corresponding formally published CC BY3.0 Unported trusted-state counting result. Nonexclusive nativev3 remains private; public selected files exclude all unselected raw ranges and comments. Full source-local Sections4.1–4.4/4.6 and AppendixA preserved, matched algorithms and specific-unitary result combined. Unused classical-sample sketch4.5 and later discussion/Followup section excluded. Original author probabilities, parameter choices and equations retained without mathematical repair or added assumptions. Original theorem/section numbering restored across omitted classicalPropositions29/30. Identity glyph uses ordinary bold1 instead of unavailable custom bbold font; no mathematical change. Bibliography identities matched to formal numbered entries. All production/validation onlyYupeng; root owns compilation.\par\medskip
\noindent\textit{Allowed standard prerequisites (editorial).}\begin{itemize}
\item Quantum query algorithms and QMA amplification; approximation theory and symmetric polynomial symmetrization. Full originalSection2 definitions/statements, source-local symmetrization proofs and guessing lemma retained; standard cited amplification is not recursively reproved
\end{itemize}\medskip
\section*{Trusted quantum samples and reflections: model and promises}
In the $\ApxCount$ problem, our goal is to decide whether a nonempty set $S\subseteq[N]$ satisfies $|S|\ge2w$ (YES) or $|S|\le w$ (NO), promised that one of these is the case. In the query model the algorithm is given a membership oracle for $S$.
QSampling $S$ means that we can approximately prepare the uniform superposition
\[|S\rangle:=\frac1{\sqrt{|S|}}\sum_{i\in S}|i\rangle.\]
To be clear, QSampling supposes that the algorithm is given trusted copies of $|S\rangle$; unlike in the QMA setting, the state need not be ``verified'' by the algorithm.
Reflections about $|S\rangle$ apply the unitary transformation
\[\mathcal R_S:=\id-2|S\rangle\langle S|,\]
which has eigenvalue $-1$ for $|S\rangle$ and eigenvalue $+1$ for all states orthogonal to $|S\rangle$. More concretely, let $U$ be the unitary that performs the map $U|0\rangle=|S\rangle$, for some canonical starting state $|0\rangle$. Since we know the circuit $U$, we can also implement $U^\dagger$, by reversing the order of all the gates and replacing all the gates with their adjoints. Then
\[\mathcal R_S=\id-2|S\rangle\langle S|=U(\id-2|0\rangle\langle0|)U^\dagger.\]
We begin by focusing on the slightly simplified setting where the algorithm is only permitted to perform membership queries, QSamples, and reflections about the state $|S\rangle$. The lower bound below even holds for the exact-size promise $|S|=w$ versus $|S|=2w$; the full unitary model is specified in original Section4.6, retained below.

\setcounter{theorem}{3}
\begin{restatable}{theorem}{main}
\label{thm:main}
Let $Q$ be a quantum algorithm that makes $T$ queries to the membership oracle for $S$,
and uses a total of $R$ copies of $|S\>$ and reflections about $|S\>$.
If $Q$ decides whether $\left\vert S\right\vert =w$\ or $\left\vert S\right\vert =2w$
with high probability, promised that one of those is the case, then either
\begin{equation}
T=\Omega\left(  \sqrt{\frac{N}{w}}\right)  \qquad \textrm{or} \qquad
R=\Omega\left(  \min\left\{  w^{1/3},\sqrt{\frac{N}{w}}\right\}  \right).
\end{equation}
\end{restatable}

\begin{restatable}{theorem}{alg}
\label{thm:alg}
There is a quantum algorithm that solves $\ApxCount$ with high probability using
$R$ copies of $|S\>$ and reflections about $|S\>$, where $R = O\left( \min \left\{ w^{1/3} , \sqrt{\frac{N}{w}} \right\} \right)$.
\end{restatable}

\setcounter{section}{1}
\setcounter{theorem}{5}
\setcounter{equation}{4}
\section{Preliminaries}

\label{sec:prelim}

In this section we introduce some definitions and known facts about
polynomials and complexity classes.

%
\subsection{Approximation theory}
\label{s:approxprelim}
We will use several results from approximation theory,\ each of which has
previously been used (in some form) in other applications of the polynomial
method to quantum lower bounds. \ We start with the basic inequality of A.A.
Markov \cite{markov1890question}.

\begin{lemma}[Markov]
\label{markovlem}Let $p$\ be a real polynomial, and suppose that
\begin{equation}
\max_{x,y\in\left[ a,b\right] }\left\vert p\left( x\right) -p\left( y\right)
\right\vert \leq H.
\end{equation}
Then for all $x\in\left[ a,b\right] $, we have
\begin{equation}
\left\vert p^{\prime}\left( x\right) \right\vert \leq\frac{H}{b-a}\deg\left(
p\right) ^{2},
\end{equation}
where $p^{\prime}(x)$ is the derivative of $p$ at $x$.
\end{lemma}

We'll also need a bound that was explicitly stated by Paturi \cite{paturi},
and which amounts to the fact that, among all degree-$d$
polynomials that are bounded within a given range, the Chebyshev polynomials
have the fastest growth outside that range.

\begin{lemma}[Paturi]
\label{paturilem}Let $p$\ be a real polynomial, and suppose that $\left\vert
p\left( x\right) \right\vert \leq1$\ for all $\left\vert x\right\vert \leq1$
. \ Then for all $x\leq1+\mu$, we have
\begin{equation}
\left\vert p\left( x\right) \right\vert \leq\exp\left( 2\deg\left( p\right)
\sqrt{2\mu+\mu^{2}}\right) .
\end{equation}
\end{lemma}

We now state a useful corollary of \Cref{paturilem}, which says (in effect)
that slightly shrinking the domain of a low-degree real polynomial can only
modestly shrink its range.

\begin{corollary}
\label{paturicor}Let $p$\ be a real polynomial of degree $d$, and suppose
that
\begin{equation}
\max_{x,y\in\left[ a,b\right] }\left\vert p\left( x\right) -p\left( y\right)
\right\vert \geq H.
\end{equation}
Let $\varepsilon\leq\frac{1}{100d^{2}}$\ and $a^{\prime}:=a+\varepsilon
\left( b-a\right) $. \ Then
\begin{equation}
\max_{x,y\in\left[ a^{\prime},b\right] }\left\vert p\left( x\right) -p\left(
y\right) \right\vert \geq\frac{H}{2}.
\end{equation}
\end{corollary}

\begin{proof}
Suppose by contradiction that
\begin{equation}
\left\vert p\left(  x\right)  -p\left(  y\right)  \right\vert <\frac{H}{2}
\end{equation}
for all $x,y\in\left[  a^{\prime},b\right]  $. \ By affine shifts, we can
assume without loss of generality that $\left\vert p\left(  x\right)
\right\vert <\frac{H}{4}$\ for all $x\in\left[  a^{\prime},b\right]  $. \ Then
by \Cref{paturilem}, for all $x\in\left[  a,b\right]  $\ we have
\begin{equation}
\left\vert p\left(  x\right)  \right\vert <\frac{H}{4}\cdot\exp\left(
2d\sqrt{2\left(  \frac{1}{1-\varepsilon}-1\right)  +\left(  \frac
{1}{1-\varepsilon}-1\right)  ^{2}}\right)  \leq\frac{H}{2}.
\end{equation}
But this violates the hypothesis.
\end{proof}

We will also need a bound that relates the range of a low-degree polynomial
on a discrete set of points to its range on a continuous interval. \ The
following lemma generalizes a result due to Ehlich and Zeller \cite{ez} and
Rivlin and Cheney \cite{rc}, who were interested only in the case where the
discrete points are evenly spaced.

\begin{lemma}
\label{ezrclem}Let $p$\ be a real polynomial of degree at most $\sqrt{k}$,
and let $0=z_{1}<\cdots<z_{M}=k$\ be a list of points such that $
z_{i+1}-z_{i}\leq1$\ for all $i$ (the simplest example being the integers $
0,\ldots,k$). \ Suppose that
\begin{equation}
\max_{x,y\in\left[ 0,k\right] }\left\vert p\left( x\right) -p\left( y\right)
\right\vert \geq H.
\end{equation}
Then
\begin{equation}
\max_{i,j}\left\vert p\left( z_{i}\right) -p\left( z_{j}\right) \right\vert
\geq\frac{H}{2}.
\end{equation}
\end{lemma}

\begin{proof}
Suppose by contradiction that
\begin{equation}
\left\vert p\left(  z_{i}\right)  -p\left(  z_{j}\right)  \right\vert
<\frac{H}{2}
\end{equation}
for all $i,j$. \ By affine shifts, we can assume without loss of generality
that $\left\vert p\left(  z_{i}\right)  \right\vert <\frac{H}{4}$\ for all
$i$. \ Let
\begin{equation}
c:=\max_{x\in\left[  0,k\right]  }\frac{\left\vert p\left(  x\right)
\right\vert }{H/4}.
\end{equation}
If $c\leq1$, then the hypothesis clearly fails, so assume $c>1$. \ Suppose
that the maximum, $\left\vert p\left(  x\right)  \right\vert =\frac{cH}{4}$,
is achieved between $z_{i}$\ and $z_{i+1}$. \ Then by basic calculus, there
exists an $x^{\ast}\in\left[  z_{i},z_{i+1}\right]  $\ such that
\begin{equation}
\left\vert p^{\prime}\left(  x^{\ast}\right)  \right\vert >\frac{2\left(
c-1\right)  }{z_{i+1}-z_{i}}\cdot\frac{H}{4}\geq\frac{\left(  c-1\right)
H}{2}.
\end{equation}
So by \Cref{markovlem},
\begin{equation}
\frac{\left(  c-1\right)  H}{2}<\frac{cH/4}{k}\deg\left(  p\right)  ^{2}.
\end{equation}
Solving for $c$, we find
\begin{equation}
c<\frac{2k}{2k-\deg\left(  p\right)  ^{2}}\leq2.
\end{equation}
But if $c<2$, then $\max_{x\in\left[  0,k\right]  }\left\vert p\left(
x\right)  \right\vert <\frac{H}{2}$, which violates the hypothesis.
\end{proof}

We also use a related inequality due to Coppersmith and Rivlin~\cite
{coppersmith-rivlin} that bounds a polynomial on a continuous interval in
terms of a bound on a discrete set of points, but now with the weaker
assumption that the degree is at most $k$, rather than $\sqrt{k}$. \ This
gives a substantially weaker bound.

\begin{lemma}[Coppersmith and Rivlin]
\label{lem:coppersmith_rivlin} Let $p$ be a real polynomial of degree at
most $k$, and suppose that $\left\vert p(x)\right\vert \leq 1$ for all
integers $x\in \{0,1,\ldots ,k\}$. \ Then there exist universal constants $
a,b$ such that for all $x\in \lbrack 0,k]$, we have
\begin{equation}
\left\vert p(x)\right\vert \leq a\cdot \exp \left( b\,\deg (p)^{2}/k\right) .
\end{equation}
\end{lemma}


%
\subsection{Symmetric polynomials}

\para{Univariate symmetrizations.}
Our starting point is the well-known \textit{symmetrization lemma} of Minsky
and Papert \cite{mp} (see also Beals et al.\ \cite{bbcmw}\ for its
application to quantum query complexity), by which we can often reduce
questions about multivariate polynomials to questions about univariate ones.

\begin{lemma}[Minsky--Papert symmetrization]
\label{symlem} Let $p:\left\{ 0,1\right\} ^{N}\rightarrow\mathbb{R}$\ be a
real multilinear polynomial of degree $d$, and let $q:\{0,1,\ldots,N\}\to
\mathbb{R}$ be defined as
\begin{equation}
q\left( k\right) :=\E_{\left\vert X\right\vert =k}\left[ p\left( X\right) 
\right] .
\end{equation}
Then $q$ can be written as a real polynomial in $k$ of degree at most $d$.
\end{lemma}

We now introduce a different, lesser known
notion of symmetrization, 
which we call the \emph{erase-all-subscripts} symmetrization
for reasons to be explained shortly. This symmetrization previously appeared in \cite{shi} under the name ``linearization,'' and it is also equivalent to the noise operator used in analysis of Boolean functions \cite[Definition 2.46]{odonnell}.

\begin{lemma}[Erase-all-subscripts
symmetrization]
Let $p:\left\{ 0,1\right\} ^{N}\rightarrow\mathbb{R}$\ be a
real multilinear polynomial of degree $d$,
and for
any real number $k \in [0, 1]$,
let $M_{k}$ denote
the distribution over $\{0, 1\}^N$, wherein each coordinate is selected
independently to be 1 with probability $k$.
Let $q:[0, 1] \to
\mathbb{R}$ be defined as
\begin{equation}
q\left( k\right) :=\E_{X \sim M_{k}}\left[ p\left( X\right) 
\right] .
\end{equation}
Then $q$ can be written as a real polynomial in $k$ of degree at most $d$.
\label{lem:eraseallsubscripts}
\end{lemma}
\begin{proof} (see, for example, \cite[Proof of Theorem 3]{STT12}).
Given the multivariate polynomial expansion of $p$,
we can obtain $q$ easily just by ``erasing all the subscripts in each variable''. For example, if
$p(x_1, x_2, x_3)= 2 x_1 x_2 + x_2 x_3 + x_2$, we replace every $x_i$ with $k$ to obtain
$q(k)=2k\cdot k+ k\cdot k + k=3k^2 + k$.
This follows from linearity of expectation along with the fact that $M_k$
is defined to be the
product distribution wherein
each coordinate has expected value $k$.
\end{proof}

We highlight the following key difference between Minsky--Papert symmetrization
and the erase-all-subscripts symmetrization. Let $p:\left\{ 0,1\right\} ^{N}\rightarrow [0, 1]$
be a real multivariate polynomial whose evaluations at Boolean inputs are in
$[0, 1]$, i.e., for all $x\in\{0,1\}^n$, we have $p(x)\in[0,1]$.
If $q$ is the erase-all-subscripts symmetrization of $p$, 
then $q$ takes values in $[0,1]$ at all \emph{real-valued}
inputs in $[0, 1]$: $q(k) \in [0,1]$ for all $k \in [0,1]$. 
If $q$ is the Minsky--Papert symmetrization of $p$,
then it is only guaranteed to take values in $[0,1]$ at 
\emph{integer-valued} inputs in $[0, N]$, i.e., $q(k)\in [0,1]$ 
is only guaranteed to hold at $k\in\{0,1,\ldots,N\}$.
This is the main reason we use erase-all-subscripts symmetrization
in this work.

\para{Bivariate symmetrizations.}
In this paper, it will be convenient
to consider bivariate versions of both
Minsky--Papert and erase-all-subscripts
symmetrization, and their applications to oracle separations. To this end,
define $X\in \left\{ 0,1\right\} ^{N}$, the \textquotedblleft characteristic
string\textquotedblright\ of the set $S\subseteq \left[ N\right] $, by $
x_{i}=1$\ if $i\in S$\ and $x_{i}=0$\ otherwise. \ Let $\mathcal{O}_{S}$
denote the unitary that performs a membership query to $S$, defined as
\begin{equation}
\mathcal{O}_{S}\left\vert i\right\rangle \left\vert b\right\rangle
=(1-2bx_{i})\left\vert i\right\rangle \left\vert b\right\rangle
\end{equation}
for any index $i\in \lbrack N]$ and bit $b\in \{0,1\}$.

Because we study oracle intersection problems, it is often convenient to
think of an algorithm as having access to \textit{two} oracles, wherein the
first bit in the oracle register selects the choice of oracle. \ As a
consequence, we need a slight generalization of a now well-established fact
in quantum complexity: that the acceptance probability of a quantum
algorithm with an oracle can be expressed as a polynomial in the bits of the
oracle string.

\begin{lemma}[Symmetrization with two oracles]
\label{lem:symmetrization} Let $Q^{\mathcal{O}_{S_{0}},\mathcal{O}_{S_{1}}}$
be a quantum algorithm that makes $T$ queries to a pair of membership
oracles for sets $S_{0},S_{1}\subseteq \lbrack N]$. \ Let $D_{\mu }$ denote
the distribution over subsets of $[N]$ wherein each element is selected
independently with probability $\frac{\mu }{N}$. \ Then there exist
bivariate real polynomials $q(s,t)$ and $p(x,y)$ of degree at most $2T$
satisfying:
\begin{align*}
\textrm{for all real numbers } s, t \in [0, N],& \quad
q(s,t)=\E_{\substack{ S_{0}\sim D_{s}, \\ S_{1}\sim D_{t}}}\left[ \Pr [Q^{
\mathcal{O}_{S_{0}},\mathcal{O}_{S_{1}}}\text{ accepts}]\right], \textrm{ and} \\
\textrm{for all integers } x, y \in \{0, 1, \dots, N\},& \quad
p(x,y)=\E_{\substack{ |S_{0}|=x, \\ |S_{1}|=y}}\left[ \Pr [Q^{\mathcal{O}
_{S_{0}},\mathcal{O}_{S_{1}}}\text{ accepts}]\right].
\end{align*}
\end{lemma}

\begin{proof}
Take $X = X_0|X_1$ to be the concatenation of the characteristic strings of the two oracles, and let $S \subseteq [2N]$ be such that $X$ is the characteristic string of $S$. Then, Lemma 4.2 of Beals et al.\ \cite{bbcmw} tells us that there is a real multilinear polynomial $r(X)$ of degree at most $2T$ in the bits of $X$ such that $r(X) = \Pr[ Q^{\mathcal{O}_S} \text{ accepts}]$.

Observe that $r$ has a meaningful probabilistic interpretation over arbitrary inputs in $[0, 1]$. A vector $X \in [0, 1]^{2N}$ of probabilities corresponds to a distribution over $\{0,1\}^{2N}$ wherein each bit is chosen from a Bernoulli distribution with the corresponding probability. Because $r$ is multilinear, $r$ in fact computes the expectation of the acceptance probability over this distribution. In particular, the polynomial
\begin{equation}
q(s, t) = r\bigg(\underbrace{\frac{s}{N},\ldots,\frac{s}{N}}_{N \text{ times}}, \underbrace{\frac{t}{N},\ldots,\frac{t}{N}}_{N \text{ times}}\bigg) = \E_{\substack{S_0 \sim D_s,\\S_1 \sim D_t}}\left[\Pr[ Q^{\mathcal{O}_{S_0},\mathcal{O}_{S_1}} \text{ accepts}] \right]    
\end{equation}
corresponds to selecting $S_0 \sim D_s$ and $S_1 \sim D_t$. The total degree of $q$ is obviously
at most the degree of $r$, by the same reasoning as in the proof of \Cref{lem:eraseallsubscripts}.

To construct $p$, we apply the symmetrization lemma of Minsky and Papert \cite{mp} to symmetrize $r$, first with respect to $X_0$, then with respect to $X_1$:
\begin{equation}
    p_0(x, X_1) = \E_{|S_0|=x} r(X_0,X_1) =  \E_{|S_0|=x}\left[\Pr[ Q^{\mathcal{O}_{S_0},\mathcal{O}_{S_1}} \text{ accepts}] \right]
\end{equation}
\begin{equation}
    p(x, y) = \E_{|S_1|=y} p_0(x,X_1) = \E_{\substack{|S_0|=x,\\|S_1|=y}}\left[\Pr[ Q^{\mathcal{O}_{S_0},\mathcal{O}_{S_1}} \text{ accepts}] \right] 
\end{equation}  
The degree of $p$ is at most the degree of $r$, due to \Cref{symlem}.
\end{proof}


We remark that, as a consequence of their definitions in \Cref{lem:symmetrization}, $p$ and $q$ satisfy:
\begin{equation}
q(s, t) = \E \left[ p(X, Y) \right],
\end{equation}
where $X$ and $Y$ are drawn from $N$-trial binomial distributions with means $s$ and $t$, respectively.

\para{Symmetric Laurent polynomials.}
Finally, we state a useful fact about Laurent polynomials:

\begin{lemma}[Symmetric Laurent polynomials]
\label{lem:symmetric} Let $\ell (x)$ be a real Laurent polynomial of positive
and negative degree $d$ that satisfies $\ell (x)=\ell (1/x)$. \ Then there exists a (ordinary) real
polynomial $q$ of degree $d$ such that $\ell (x)=q(x+1/x)$.
\end{lemma}

\begin{proof}
$\ell(x) = \ell(1/x)$ implies that the coefficients of the $x^i$ and $x^{-i}$ terms are equal for all $i$, as otherwise $\ell(x) - \ell(1/x)$ would not equal the zero polynomial. Thus, we may write $\ell(x) = \sum_{i=0}^d a_i \cdot (x^i + x^{-i})$ for some coefficients $a_i$. So, it suffices to show that $x^i + x^{-i}$ can be expressed as a polynomial in $x + 1/x$ for all $0 \le i \le d$.

We prove by induction on $i$. The case $i = 0$ corresponds to constant polynomials. For $i > 0$, by the binomial theorem, observe that $(x + 1/x)^i = x^i + x^{-i} + r(x)$ where $r$ is a degree $i - 1$ real Laurent polynomial satisfying $r(x) = r(1/x)$. By the induction assumption, $r$ can be expressed as a polynomial in $x + 1/x$, so we have $x^i + x^{-i} = (x + 1/x)^i - r(x)$ is expressed as a polynomial in $x + 1/x$.
\end{proof}

%

\setcounter{section}{3}
\setcounter{theorem}{20}
\section{Approximate counting with quantum samples and reflections}

\label{sec:samplesreflections}

%
\subsection{The Laurent polynomial method}

\label{RESULT}\label{sec:lower}\label{sec:Laurent}

By using Minsky--Papert symmetrization (\Cref{symlem}), we now prove the key
fact that relates quantum algorithms, of the type we're considering, to real
Laurent polynomials in one variable. \ The following lemma generalizes the
connection between quantum algorithms and real polynomials established by
Beals et al.\ \cite{bbcmw}.

\begin{lemma}
\label{laurentlem} Let $Q$ be a quantum algorithm that makes $T$ queries to $
\mathcal{O}_{S}$, uses $R_1$ copies of $\left\vert S\right\rangle $, and
makes $R_2$ uses of the unitary $\mathcal{R}_S$. \ Let $R:= R_1 + 2R_2$. 
For $k\in\{1,\ldots,N\}$, let
\begin{equation}
q\left( k\right) :=\E_{\left\vert S\right\vert =k}\left[ \Pr\left[ Q^{
\mathcal{O}_{S},\mathcal{R}_S}\left( \left\vert S\right\rangle ^{\otimes
R_1}\right) \text{ accepts}\right] \right] .
\end{equation}
Then $q$ can be written a univariate Laurent polynomial, with maximum 
exponent at most $2T+R$\ and minimum exponent at least $-R$.
\end{lemma}

\begin{proof}
Let $|\psi_\mathrm{initial}\>$ denote the initial state of the algorithm,
which we can write as
\begin{align*}
|\psi_\mathrm{initial}\> &=
\left\vert S\right\rangle ^{\otimes R_1}
=\left(\frac{1}{\sqrt{|S|}} \sum_{i\in S}|i\>\right)^{\otimes R_1}
=\left(\frac{1}{\sqrt{|S|}} \sum_{i\in [N]}x_i|i\>\right)^{\otimes R_1}\\
&=\frac{1}{\left\vert S\right\vert ^{R_1/2}}\sum_{i_{1},\ldots,i_{R_1}\in\left[
N\right]  }x_{i_{1}}\cdots x_{i_{R_1}}\left\vert i_{1},\ldots,i_{R_1}\right\rangle.
\end{align*}
Thus, each amplitude is a complex multilinear polynomial in $X=\left(
x_{1},\ldots,x_{N}\right)  $ of degree $R_1$, divided by $\left\vert
S\right\vert ^{R_1/2}$.

Throughout the algorithm, each amplitude will remain a complex multilinear
polynomial in $X$ divided by some power of $|S|$.
Since $x_{i}^{2}=x_{i}$\ for all $i$, we can always
maintain multilinearity without loss of generality.

Like Beals et al.\ \cite{bbcmw}, we now consider how the polynomial degree of each amplitude and the power of $|S|$ in the denominator change as the algorithm progresses.
We have to handle 3 different kinds of unitaries that the quantum circuit may use:
the membership query oracle $\mathcal{O}_S$, unitaries independent of the input,
and the reflection unitary $\R_S$.

The first two cases are handled as in Beals et al.\
Since $\mathcal{O}_S$ is a unitary whose entries are degree-1 polynomials
in $X$, each use of this unitary increases a particular
amplitude's degree as a polynomial by $1$ and does not change the power of $|S|$ in the denominator. \
Second, input-independent unitary transformations only take linear
combinations of existing polynomials and hence do not increase the degree of the amplitudes or the power of $|S|$ in the denominator.
Finally, we consider the reflection unitary $\R_S = \id - 2|S\>\<S|$.
The $(i,j)^\mathrm{th}$ entry of this operator is $\delta_{ij} - \frac{2x_ix_j}{|S|} = \frac{\delta_{ij}|S| - 2x_ix_j}{|S|}$, where $\delta_{ij}$ is the Kronecker delta function.
Since $|S| = \sum_i x_i$, this is a degree-2 polynomial divided by $|S|$.
Hence applying this unitary will increase the degree of the amplitudes by $2$ and increase
the power of $|S|$ in the denominator by $1$.

In conclusion, we start with each amplitude being a polynomial of degree $R_1$ divided by $|S|^{R_1/2}$. $T$ queries to the membership oracle will increase the degree of each amplitude by at most $T$ and leave the power of $|S|$ in the denominator unchanged. $R_2$ uses of the reflection unitary will increase the degree by at most $2R_2$ and the power of $|S|$ in the denominator by $R_2$. It follows that $Q$'s final state has the form
\begin{equation}
\left\vert \psi_\mathrm{final}\right\rangle =\sum_z \alpha_{z}\left(  X\right)
\left\vert z\right\rangle ,
\end{equation}
where each $\alpha_{z}\left(  X\right)  $\ is a complex multilinear polynomial
in $X$ of degree at most $R_1+2R_2+T = R+T$, divided by $\left\vert S\right\vert
^{R_1/2+R_2} = |S|^{R/2}$. \ Since $X$ itself is real-valued, it follows that the real and
imaginary parts of $\alpha_{z}\left(  X\right)  $, considered individually,
are real multilinear polynomials in $X$\ of degree at most $R+T$\ divided by
$\left\vert S\right\vert ^{R/2}$.

Hence, if we let
\begin{equation}
p\left(  X\right)  :=\Pr\left[  Q^{\mathcal{O}_{S},\R_S}\left(  \left\vert
S\right\rangle ^{\otimes R_1}\right)  \text{ accepts}\right]  ,
\end{equation}
then
\begin{equation}
p\left(  X\right)  =\sum_{\text{accepting }z}\left\vert \alpha_{z}\left(
X\right)  \right\vert ^{2}=\sum_{\text{accepting }z}\left(  \operatorname{Re}
^{2}\alpha_{z}\left(  X\right)  +\operatorname{Im}^{2}\alpha_{z}\left(
X\right)  \right)
\end{equation}
is a real multilinear polynomial in $X$ of degree at most $2\left(
R+T\right)  $, divided through (in every monomial) by $\left\vert S\right\vert
^{R}=\left\vert X\right\vert ^{R}$.

Now consider
\begin{equation}
q\left(  k\right)  :=\E_{\left\vert X\right\vert =k}\left[
p\left(  X\right)  \right]  .
\end{equation}
By \Cref{symlem}, this is a real univariate polynomial in $\left\vert
X\right\vert $ of degree at most $2\left(  R+T\right)  $, divided through (in
every monomial) by $\left\vert S\right\vert ^{R}=\left\vert X\right\vert ^{R}
$. \ Or said another way, it's a real Laurent polynomial in $\left\vert
X\right\vert $, with maximum exponent at most $R+2T$\ and minimum exponent at
least $-R$.
\end{proof}

%
\subsection{Upper bounds}
\label{UPPER}\label{sec:upper}

Before proving our lower bounds on the degree of 
Laurent polynomials approximating $\mathsf{ApxCount}_{N,w}$, we
establish some simpler \emph{upper bounds}.
We show upper bounds on Laurent polynomial degree and in the 
queries, samples, and reflections model.

\para{Laurent polynomial degree of approximate counting.}
We now describe a \emph{purely negative} degree Laurent polynomial of degree $O(w^{1/3})$ for
approximate counting. \ This upper bound will serve as an important source
of intuition when we prove the (matching) lower bound of \Cref{thm:main}
(see \Cref{s:nextsubsection}). We are thankful to user \textquotedblleft
fedja\textquotedblright\ on MathOverflow for describing this construction.
\footnote{
See 
\url{https://mathoverflow.net/questions/302113/real-polynomial-bounded-at-inverse-integer-points}
}

\begin{lemma}[fedja]
\label{fedjatight} For all $w$, there is a real polynomial $p$ of degree $O\left(w^{1/3}\right)$ such that:
\begin{enumerate}
\item $0 \le p(1/k) \le \frac{1}{3}$ for all $k \in [w]$.
\item $\frac{2}{3} \le p(1/k) \le 1$ for all integers $k \ge 2w$.
\item $0 \le p(1/k) \le 1$ for all $k \in \{w+1,w+2,\ldots,2w-1\}$.
\end{enumerate}
\end{lemma}

\begin{proof}
Assuming for simplicity that $w$ is a perfect cube, consider
\begin{equation}
u\left(  x\right)  :=\left(  1-x\right)  \left(  1-2x\right)  \cdots\left(
1-w^{1/3}x\right)  .
\end{equation}
Notice that $\deg\left(  u\right)  =w^{1/3}$\ and $u\left(  \frac{1}
{k}\right)  =0$ for all $k\in\left[  w^{1/3}\right]  $. \ Furthermore, we have
$u\left(  x\right) \in [0,1]$\ for all $x\in\left[
0,\frac{1}{w^{1/3}}\right]  $, and also $u\left(  x\right)  \in\left[
1-O\left(  \frac{1}{w^{1/3}}\right)  ,1\right]  $\ for all $x\in\left[
0,\frac{1}{w}\right]  $. \ Now, let $v$\ be the Chebyshev polynomial of degree
$w^{1/3}$, affinely adjusted so that $v\left(  x\right)
\in [0,1]$\ for all $x\in\left[  0,\frac{1}{w^{1/3}}\right]  $ (rather
than in $[-1,1]$ for all all $\left\vert x\right\vert \leq1$), and with a large jump between
$\frac{1}{2w}$\ and $\frac{1}{w}$. \ Then the product, $p(x):=u\left(  x\right)  v\left(  x\right)  $, has degree $2w^{1/3}$\ and
satisfies all the requirements, except possibly that the constants $\frac{1}{3}$ and $\frac{2}{3}$ in the first two requirements may be off. Composing with a constant degree polynomial corrects this, and gives a polynomial of degree $O(w^{1/3})$ that satisfies all three requirements.
\end{proof}

Interestingly, if we restrict our attention to purely negative degree
Laurent polynomials, then a matching lower bound is not too hard to show. \
In the same MathOverflow post, user fedja also proves the following, which can also
be shown using earlier work of Zhandry \cite[Proof of Theorem 7.3]{zhandry}):

\begin{lemma}
\label{fedjalem}Let $p$\ be a real polynomial, and suppose that $\left\vert
p\left( 1/k\right) \right\vert \leq1$ for all $k\in\left[ 2w\right] $, and
that $p\left( \frac{1}{w}\right) \leq\frac{1}{3}$\ while $p\left( \frac{1}{2w
}\right) \geq\frac{2}{3}$. \ Then $\deg\left( p\right) =\Omega\left(
w^{1/3}\right) $.
\end{lemma}

\Cref{s:explosion} and \Cref{s:dualpoly} below take the considerable step of
extending \Cref{fedjalem} from purely negative degree Laurent
polynomials to general Laurent polynomials.


\para{Upper bounds in the queries, samples, and reflections model.} 
Although we showed that there is a purely negative degree Laurent polynomial of degree
$O(w^{1/3})$ for $\ApxCount$, this does not imply the existence of a quantum algorithm
in the queries, samples, and reflections model with similar complexity.

We now show that our lower bounds in the queries, samples, and reflections model (in \Cref{thm:main}) are tight (up to constants). This is \Cref{thm:alg} in the introduction, restated here for convenience:

\alg*

\begin{proof}
We describe two quantum algorithms for this problem with the two stated complexities.

The first algorithm uses $O(w^{1/3})$ samples and reflections. This algorithm is reminiscent of the original collision finding algorithm of Brassard, H{\o}yer, and Tapp~\cite{BHT98}. We first use $O(w^{1/3})$ copies of $|S\>$ to learn a set $M\subset S$ of size $w^{1/3}$ by simply measuring copies of $|S\>$ in the computational basis. Now we know that the ratio $|S|/|M|$ is either $w^{2/3}$ or $2w^{2/3}$. Now consider running Grover's algorithm on the set $S$ where the elements in $M$ are considered the ``marked'' elements. Grover's algorithm alternates reflections about the uniform superposition over the set being searched, $S$, with an operator that reflects about the marked elements in $M$. The first reflection is simply $\R_S$, which we have access to. The second unitary can be constructed since we have an explicit description of the set $M$. Now Grover's algorithm can be used to distinguish whether the fraction of marked elements is $1/w^{2/3}$ or half of that, and the cost will be $O(w^{1/3})$.

The second algorithm uses $O(\sqrt{{N}/{w}})$ reflections only and no copies of $|S\>$. 
Consider running the standard approximate counting algorithm~\cite{BHMT02} that uses membership queries to $S$ and distinguishes $|S|\leq w$ from $|S|\geq 2w$ using $O(\sqrt{N/w})$ membership queries. 
Observe that this algorithm starts with the state $|\psi\> =\frac{1}{
\sqrt{N}}\left( \left\vert 1\right\rangle +\cdots +\left\vert N\right\rangle
\right)$, which is in $\textrm{span}\{|S\>,|\bar{S}\>\}$, and only uses reflections about $|\psi\> $ and membership queries to $|S\>$ in the form of a unitary that maps $|i\>$ to $-|i\>$ when $i \in S$.
This means the state of the algorithm remains in $\textrm{span}\{|S\>,|\bar{S}\>\}$
at all times. 
Within this subspace, a membership query to $S$ is the same as a reflection about $|S\>$.
Hence we can replace membership queries with the reflection operator to get an approximate counting algorithm that only uses $O(\sqrt{N/w})$ reflections and no copies of $|S\>$.
\end{proof}

Note that both the algorithms presented above generalize to the situation
where we want to distinguish $|S|=w$ from $|S|=(1+\varepsilon )w$. \ For the
first algorithm, we now pick a subset $M$ of size $w^{1/3}/\varepsilon ^{2/3}
$. Now we want to $(1+\varepsilon )$-approximate the fraction of marked
elements, which is either $1/(w\varepsilon )^{2/3}$ or $(1+\varepsilon )^{-1}
$ times that. This can be done with approximate counting~\cite[Theorem 15]
{BHMT02}, and the cost will be $O\left( \frac{1}{\varepsilon }(w\varepsilon
)^{1/3}\right) =O\left( \frac{w^{1/3}}{\varepsilon ^{2/3}}\right) $.
The second algorithm is simpler to generalize, since we simply plug in
the query complexity of $\varepsilon$-approximate counting, which is 
$O\Bigl( \frac{1}{\varepsilon }\sqrt{\frac{N}{w}}\Bigr) $.

%
\subsection{Lower bound using the explosion argument}
\label{s:explosion}\label{sec:explosion}

We now show a weaker version of \Cref{thm:main} using the explosion argument
described in the introduction. \ The difference between the following
theorem and \Cref{thm:main} is the exponent of $w$ in the lower bound.

\begin{theorem}
\label{thm:explosion}Let $Q$ be a quantum algorithm that makes $T$ queries
to the membership oracle for $S$, and uses a total of $R$ copies of $|S\>$
and reflections about $|S\>$. \ If $Q$ decides whether $\left\vert
S\right\vert =w$\ or $\left\vert S\right\vert =2w$ with success probability
at least $2/3$, promised that one of those is the case, then either
\begin{equation}
T=\Omega \left( \sqrt{\frac{N}{w}}\right) \qquad \text{or}\qquad R=\Omega
\left( \min \left\{ w^{1/4},\sqrt{\frac{N}{w}}\right\} \right) .
\end{equation}
\end{theorem}

\begin{proof}
Since we neglect multiplicative constants in our lower bounds, let us allow the algorithm to use up to $R$ copies of $|S\>$ and $R$ uses of $\R_S$.
Let
\begin{equation}
q\left(  k\right)  :=\E_{\left\vert S\right\vert =k}\left[
\Pr\left[  Q^{\mathcal{O}_{S},\R_S}\left(  \left\vert S\right\rangle ^{\otimes
R}\right)  \text{ accepts}\right]  \right]  .
\end{equation}
Then by \Cref{laurentlem}, we can write $q$ as a Laurent polynomial:
\begin{equation}
q\left(  k\right)  =u\left(  k\right)  +v\left(  1/k\right)  ,
\end{equation}
where $u$\ is a real polynomial in $k$ with $\deg\left(  u\right)  = O(T+R)$,
and $v$\ is a real polynomial in $1/k$\ with $\deg\left(  v\right) = O(R)$.
\ So to prove the theorem, it suffices to show that either $\deg\left(
u\right)  =\Omega\left(  \sqrt{\frac{N}{w}}\right)  $, or else $\deg\left(
v\right)  =\Omega\left(  w^{1/4}\right)  $. \ To do so, we'll assume that
$\deg\left(  u\right)  =o\left(  \sqrt{\frac{N}{w}}\right)  $\ and
$\deg\left(  v\right)  =o\left(  w^{1/4}\right)  $, and derive a contradiction.

Our high-level strategy is as follows: we'll observe that, if approximate
counting is being successfully solved, then either $u$ or $v$ must attain a
large first derivative somewhere in its domain. \ By the approximation theory
lemmas that we proved in \Cref{s:approxprelim}, this will force that polynomial to have a large
range---even on a subset of integer (or inverse-integer) points. \ But the
sum, $u\left(  k\right)  +v\left(  1/k\right)  $, is bounded in $\left[
0,1\right]  $\ for all $k\in\left[  N\right]  $. \ So if one polynomial has a
large range, then the other does too. \ But this forces the \textit{other}
polynomial to have a large derivative somewhere in its domain, and therefore
(by approximation theory) to have an even larger range, forcing the first
polynomial to have an even larger range to compensate, and so on. \ As long as
$\deg\left(  u\right)  $\ and $\deg\left(  v\right)  $ are both small enough,
this endless switching will force both $u$ and $v$ to attain
\textit{unboundedly }large values---with the fact that one polynomial is in
$k$, and the other is in $1/k$, crucial to achieving the desired
\textquotedblleft explosion.\textquotedblright\ \ Since $u$ and $v$ are
polynomials on compact sets, such unbounded growth is an obvious absurdity,
and this will give us the desired contradiction.

In more detail, we will study the following quantities.
\begin{equation}
\begin{tabular}
[c]{ll}
$G_{u}:=\max_{x,y\in\left[  \sqrt{w},2w\right]  }\left\vert u\left(  x\right)
-u\left(  y\right)  \right\vert ~\ \ \ \ \ \ $ & $G_{v}:=\max_{x,y\in\left[
\frac{1}{N},\frac{1}{w}\right]  }\left\vert v\left(  x\right)  -v\left(
y\right)  \right\vert $\\
$\Delta_{u}:=\max_{x\in\left[  \sqrt{w},2w\right]  }\left\vert u^{\prime
}\left(  x\right)  \right\vert $ & $\Delta_{v}:=\max_{x\in\left[  \frac{1}
{N},\frac{1}{w}\right]  }\left\vert v^{\prime}\left(  x\right)  \right\vert
$\\
$H_{u}:=\max_{x,y\in\left[  \sqrt{w},N\right]  }\left\vert u\left(  x\right)
-u\left(  y\right)  \right\vert $ & $H_{v}:=\max_{x,y\in\left[  \frac{1}
{N},\frac{1}{\sqrt{w}}\right]  }\left\vert v\left(  x\right)  -v\left(
y\right)  \right\vert $\\
$I_{u}:=\max_{x,y\in\left[  w,N\right]  }\left\vert u\left(  x\right)
-u\left(  y\right)  \right\vert $ & $I_{v}:=\max_{x,y\in\left[  \frac{1}
{2w},\frac{1}{\sqrt{w}}\right]  }\left\vert v\left(  x\right)  -v\left(
y\right)  \right\vert $\\
$L_{u}:=\max_{x,y\in\left\{  w,\ldots,N\right\}  }\left\vert u\left(
x\right)  -u\left(  y\right)  \right\vert $ & $L_{v}:=\max_{x,y\in\left\{
\sqrt{w},\ldots,2w\right\}  }\left\vert v\left(  \frac{1}{x}\right)  -v\left(
\frac{1}{y}\right)  \right\vert $
\end{tabular}
\end{equation}


We have $0\leq q\left(  k\right)  \leq1$\ for all $k\in\left[  N\right]  $,
since in those cases $q\left(  k\right)  $\ represents a probability. \ Since
$Q$ solves approximate counting, we also have $q\left(  w\right)  \leq\frac
{1}{3}$\ and $q\left(  2w\right)  \geq\frac{2}{3}$. \ This means in particular
that either

\begin{enumerate}
\item[(i)] $u\left(  2w\right)  -u\left(  w\right)  \geq\frac{1}{6}$, and
hence $G_{u}\geq\frac{1}{6}$, or else

\item[(ii)] $v\left(  \frac{1}{2w}\right)  -v\left(  \frac{1}{w}\right)
\geq\frac{1}{6}$, and hence $G_{v}\geq\frac{1}{6}$.
\end{enumerate}

We will show that either case leads to a contradiction.

We have the following inequalities regarding $u$:
\begin{equation}
\begin{tabular}
[c]{ll}
$G_{u}\geq L_{v}-1$ & by the boundedness of $q$\\
$\Delta_{u}\geq\frac{G_{u}}{2w}$ & by basic calculus\\
$H_{u}\geq\frac{\Delta_{u}\left(  N-\sqrt{w}\right)  }{\deg\left(  u\right)
^{2}}$ & by \Cref{markovlem}\\
$I_{u}\geq\frac{H_{u}}{2}$ & by \Cref{paturicor}\\
$L_{u}\geq\frac{I_{u}}{2}$ & by \Cref{ezrclem}
\end{tabular}
\end{equation}
Here the fourth inequality uses the fact that, setting $\varepsilon
:=\frac{\sqrt{w}}{N}$, we have $\deg\left(  u\right)  =o\left(  \frac{1}
{\sqrt{\varepsilon}}\right)  $ (thereby satisfying the hypothesis of \Cref{paturicor}), while the fifth inequality uses the fact that $\deg\left(
u\right)  =o\left(  \sqrt{N}\right)  $.

Meanwhile, we have the following inequalities regarding $v$:
\begin{equation}
\begin{tabular}
[c]{ll}
$G_{v}\geq L_{u}-1$ & by the boundedness of $q$\\
$\Delta_{v}\geq G_{v}w$ & by basic calculus\\
$H_{v}\geq\frac{\Delta_{v}\left(  \frac{1}{\sqrt{w}}-\frac{1}{N}\right)
}{\deg\left(  v\right)  ^{2}}$ & by \Cref{markovlem}\\
$I_{v}\geq\frac{H_{v}}{2}$ & by \Cref{paturicor}\\
$L_{v}\geq\frac{I_{v}}{2}$ & by \Cref{ezrclem}
\end{tabular}
\end{equation}
Here the fourth inequality uses the fact that, setting $\varepsilon
:=\frac{1/2w}{1/\sqrt{w}}=\frac{1}{2\sqrt{w}}$, we have $\deg\left(  v\right)
=o\left(  \frac{1}{\sqrt{\varepsilon}}\right)  $ (thereby satisfying the
hypothesis of \Cref{paturicor}). \ The fifth inequality uses the fact
that, if we set $V\left(  x\right)  :=v\left(  x/w\right)  $, then the
situation\ satisfies the hypothesis of \Cref{ezrclem}:\ we are interested
in the range of $V$ on the interval $\left[  \frac{1}{2},\sqrt{w}\right]  $,
compared to its range on discrete points $\frac{w}{\sqrt{w}},\frac{w}{\sqrt
{w}+1},\ldots,\frac{w}{2w}$\ that are spaced at most $1$ apart from each
other; and we also have $\deg\left(  V\right)  =\deg\left(  v\right)
=o\left(  w^{1/4}\right)  $.

All that remains is to show that, if we insert either $G_{u}\geq\frac{1}{6}
$\ or $G_{v}\geq\frac{1}{6}$ into the coupled system of inequalities above,
then we get unbounded growth and the inequalities have no solution. \ Let us
collapse the two sets of inequalities to
\begin{align*}
L_{u}  &  \geq\frac{1}{4}\frac{N-\sqrt{w}}{\deg\left(  u\right)  ^{2}}
\frac{G_{u}}{2w}=\Omega\left(  \frac{N}{w\deg\left(  u\right)  ^{2}}
G_{u}\right)  ,\\
L_{v}  &  \geq\frac{1}{4}\frac{\frac{1}{\sqrt{w}}-\frac{1}{N}}{\deg\left(
v\right)  ^{2}}G_{v}w=\Omega\left(  \frac{\sqrt{w}}{\deg\left(  v\right)
^{2}}G_{v}\right)  .
\end{align*}
Hence
\begin{align*}
G_{u}  &  \geq L_{v}-1=\Omega\left(  \frac{\sqrt{w}}{\deg\left(  v\right)
^{2}}G_{v}\right)  -1,\\
G_{v}  &  \geq L_{u}-1=\Omega\left(  \frac{N}{w\deg\left(  u\right)  ^{2}
}G_{u}\right)  -1.
\end{align*}
By the assumption that $\deg\left(  v\right)  =o\left(  w^{1/4}\right)  $\ and
$\deg\left(  u\right)  =o\left(  \sqrt{\frac{N}{w}}\right)  $, we have
$\frac{\sqrt{w}}{\deg\left(  v\right)  ^{2}}\gg1$\ and $\frac{N}{w\deg\left(
u\right)  ^{2}}\gg1$. \ Plugging in $G_{u}\geq\frac{1}{6}$\ or $G_{v}\geq
\frac{1}{6}$, this is enough to give us unbounded growth.
\end{proof}

%
\subsection{Lower bound using dual polynomials}
\label{s:dualpoly}\label{sec:dualpoly}

In this section we use the method of dual polynomials to
establish our main result, \Cref{thm:main}, restated for convenience:

\main*

Let $p(r)$ be a univariate Laurent polynomial of negative degree $D_{1}$ and
positive degree $D_{2}$. \ That is, let $p(r)$ be of the form
\begin{equation}
p(r)=a_{0}/r^{D_{1}}+a_{1}/r^{D_{1}-1}+\dots
+a_{D_{1}-1}/r+a_{D_{1}}+a_{D_{1}+1}\cdot r+\cdots +a_{D_{2}+D_{1}}\cdot
r^{D_{2}}.
\end{equation}

\Cref{thm:main} follows by combining the Laurent polynomial method (\Cref{laurentlem}) and the following
theorem.

\begin{theorem}
\label{mainthm} 
Let $\varepsilon < 1$. Suppose that $p$ has negative degree $D_1$ and positive
degree $D_2$ and satisfies the following properties.

\begin{itemize}
\item $\left\vert p(w)-1\right\vert \leq \varepsilon $

\item $\left\vert p(2w)+1\right\vert \leq \varepsilon $

\item $\left\vert p(\ell )\right\vert \leq 1+\varepsilon \text{ for all }
\ell \in \{1,2,\dots ,n\}$
\end{itemize}

Then either $D_1 \geq \Omega\left(w^{1/3}\right)$ or $D_2 \geq \Omega\left(\sqrt{N/w}
\right)$.
\end{theorem}

In fact, our proof of \Cref{mainthm} will show that the lower bound holds
even if $\left\vert p(\ell )\right\vert \leq 1+\varepsilon $ only for $\ell
\in \{w^{1/3},w^{1/3}+1,\dots ,w\}\cup \{2w,2w+1,\dots ,N\}$. \ We refer to
a Laurent polynomial $p$ satisfying the three properties of \Cref{mainthm}
as an \emph{approximation for approximate counting}.

\paragraph{Proof of \Cref{mainthm}.}

Let $p$ be any Laurent polynomial satisfying the hypothesis of \Cref{mainthm}
. \ We begin by transforming $p$ into a (standard) polynomial $q$ in a
straightforward manner. \ This transformation is captured in the following
lemma, whose proof is so simple that we omit it.

\begin{lemma}
\label{mainlem} If $p$ satisfies the properties of \Cref{mainthm}, then the
polynomial $q(r) = p(r) \cdot r^{D_1}= a_0 + a_1 r + \dots + a_{D_1 + D_2}
r^{D_1 + D_2}$ is a (standard) polynomial of degree at most $D_1 + D_2$, and
$q$ satisfies the following three properties.

\begin{itemize}
\item $\left\vert q(w)-w^{D_{1}}\right\vert \leq \varepsilon \cdot w^{D_{1}}$

\item $\left\vert q(2w)+(2w)^{D_{1}}\right\vert \leq \varepsilon \cdot
(2w)^{D_{1}}$

\item $\left\vert q(\ell )\right\vert \leq \left( 1+\varepsilon \right) \ell
^{D_{1}}\text{ for all }\ell \in \{1,2,\dots ,N\}$
\end{itemize}
\end{lemma}

We now turn to showing that, for any constant $\varepsilon < 1$, no
polynomial $q$ can satisfy the conditions of \Cref{mainlem} unless $D_1 \geq
\Omega(w^{1/3})$ or $D_2 \geq \Omega\left(\sqrt{N/w}\right)$.

Consider the following linear program. \ The variables of the linear program
are $\varepsilon $, and the $D_{2}+D_{1}+1$ coefficients of $q$.

\begin{equation}
\boxed{\begin{array}{ll} \text{minimize } & \eps \\ \mbox{such that} & \\
&|q(w) - w^{D_1}| \leq \eps \cdot w^{D_1}\\ &|q(2w) + (2w)^{D_1}| \leq \eps
\cdot (2w)^{D_1}\\ & |q(\ell)| \leq (1+\eps) \cdot \ell^{D_1} \text{ for all
} \ell \in \{1, 2, \dots, N\}\\ &\eps \geq 0 \end{array}}
\end{equation}

Standard manipulations reveal the dual.

\begin{equation} \label{duallp}
\boxed{\begin{array}{ll} \text{maximize } \phi(w) \cdot w^{D_1} - \phi(2w)
\cdot (2w)^{D_1} - \sum_{\ell \in \{1, \dots, N\}, \ell \not\in \{w, 2w\}}
|\phi(\ell)| \cdot \ell^{D_1} \\ \mbox{such that} & \\ \sum_{\ell=1}^{N}
\phi(\ell) \cdot \ell^j =0 \text{ for } j= 0, 1, 2, \dots, D_1 + D_2 &\\
\sum_{\ell=1}^N |\phi(\ell)| \cdot \ell^{D_1} = 1\\ \phi \colon \Reals \to
\mathbb{R} \end{array}}
\end{equation}

\label{s:duallp}

\medskip\medskip

\Cref{mainthm} will follow if we can exhibit a solution $\phi $ to the dual
linear program achieving value $\varepsilon >0$, for some
setting of $D_{1}\geq \Omega (w^{1/3})$ and $D_{2}\geq \Omega \left( \sqrt{
N/w}\right) $.\footnote{
We will alternatively refer to such dual solutions $\phi $ as \emph{dual
witnesses}, since they act as a witness to the fact that any low-degree
Laurent polynomial $p$ approximating the approximate counting problem must
have large error.} \ We now turn to this task.


\subsubsection{Constructing the dual solution}
\label{duallpsec}

For a set $T\subseteq \{0,1,\dots ,N\}$, define
\begin{equation}
Q_{T}(t)=\prod_{i=0,1,\dots ,N,i\not\in T}(t-i).
\end{equation}
Let $c>2$ be an integer constant that we will choose later (the bigger we choose $c$
to be, the better the objective value achieved by our final dual witness. \
But choosing a bigger $c$ will also lower the degrees $D_{1},D_{2}$ of
Laurent polynomials against which our lower bound will hold).

We now define two sets $T_{1}$ and $T_{2}$. \ The size of $T_{1}$ will be
\begin{equation}
d_{1}:=\lfloor \left( w/c\right) ^{1/3}\rfloor =\Theta \left( w^{1/3}\right)
\end{equation}
and the size of $T_{2}$ will be $d_{2}$ for
\begin{equation}
d_{2}:=\lfloor \sqrt{N/(cw)}\rfloor =\Theta \left( \sqrt{N/w}\right) .
\end{equation}
%
%

Let
\begin{equation}
T_1 = \left\{\lfloor w/(c i^2) \rfloor \colon i=1, 2, \dots, d_1\right\}
\end{equation}
and
\begin{equation}
T_2 = \left\{c \cdot i^2 \cdot w \colon i = 1, 2, \dots, d_2:=\sqrt{N/(cw)}
\right\}.
\end{equation}
Finally, define
\begin{equation}
T = \{w, 2w\} \cup T_1 \cup T_2.
\end{equation}

At last, define $\Phi \colon \{0, 1, \dots, N\} \to \mathbb{R}$ via
\begin{equation}  \label{phidef}
\Phi(t) = (-1)^t \cdot \binom{N}{t} \cdot Q_T(t).
\end{equation}

Our final dual solution $\phi $ will be a scaled version of $\Phi $. \
Specifically, $\Phi $ itself does not satisfy the second constraint of the
dual linear program, that $\sum_{\ell =1}^{N}|\Phi (\ell )|\cdot \ell
^{D_{1}}=1$. \ So letting
\begin{equation}
C=\sum_{\ell =1}^{N}|\Phi (\ell )|\cdot \ell ^{D_{1}},  \label{Cdef}
\end{equation}
our final dual witness $\phi $ will be $\Phi /C$.

\para{The sizes of $T_{1}$ and $T_{2}$.} \ Clearly,
under the above definition of $T_{2}$, $|T_{2}|=d_{2}$ as claimed above. \
It is not as immediately evident that $|T_{1}|=d_{1}$: to establish this, we
must show that for distinct $i,j\in \{1,2,\dots ,d_{1}\}$, $\lfloor
w/(ci^{2})\rfloor \neq \lfloor w/(cj^{2})\rfloor $. \ This is handled in the
following easy lemma.

\begin{lemma}
\label{lem:distinct} Let $i\neq j$ be distinct numbers in $\{1,\dots ,d_{1}\}
$ and $c>2$ be a constant. \ Then as long as $d_{1}<\left( w/c\right) ^{1/3}$
, it holds that $\lfloor w/(ci^{2})\rfloor \neq \lfloor w/(cj^{2})\rfloor $.
\end{lemma}

\begin{proof}
Assume without loss of generality that  $i> j$. Then $w/(c j^2)- w/(c i^2)$ is clearly minimized when $i=d_1$ and $j=i-1$.
For the remainder of the proof, fix $i=d_1$.
In this case,
\begin{align}
w/(c j^2) - w/(c i^2) \geq
w/\left(c (i-1)^2\right) - w/\left(c i^2\right)
= \frac{w i^2 - w \left(i-1\right)^2}{c \cdot i^2 \cdot \left(i-1\right)^2} \notag \\
= \frac{w}{c} \cdot \frac{2i-1}{i^2\left(i-1\right)^2 }
\geq \frac{w}{c} \cdot \frac{2i-1}{i^4}
\geq \frac{w}{c i^3} \geq 1. \label{eq28}
\end{align}
Here, the final inequality holds because $i^3 = d_1^3 \leq w/c$.

\Cref{eq28} implies the lemma, as two numbers whose difference is at least 1 cannot have
the same integer floor.
\end{proof}

\Cref{lem:distinct} is false for $d_1 = \omega(w^{1/3})$, highlighting on a
technical level why one cannot choose $d_1$ larger than $\Theta(w^{1/3})$
without the entire construction and analysis of $\Phi$ breaking down.

\subsubsection{Intuition: ``gluing together'' two simpler dual solutions}

Before analyzing the dual witnesses
$\Phi$ and $\phi$ constructed in \Cref{phidef} and \Cref{Cdef},
in this subsection and the next, we provide detailed intuition
for why the definitions of $\Phi$ and $\phi$ are natural,
and briefly overview their analysis. 

\label{sec:intuition}

\para{A dual witness for purely positive degree (i.e.,
approximate degree).} \ Suppose we were merely interested in showing an
approximate degree lower bound of $\Omega (\sqrt{N/w})$ for approximate
counting (i.e., a lower bound on the degree of traditional polynomials that
distinguish input $w$ from $2w$, and are bounded at all other integer inputs
in $1,\dots ,N$). \ This is equivalent to exhibiting a solution to the dual
linear program with $D_{1}=0$. \ A valid dual witness $\phi_1$ for this
simpler case is to also use \Cref{phidef}, but to set

\begin{equation}  \label{simplerT1}
T= \{w, 2w\} \cup T_2,
\end{equation}
rather than $T=\{w, 2w\} \cup T_1 \cup T_2$.

We will explain intuition for why \Cref{simplerT1} is a valid dual
solution for the approximate degree of approximate counting in the next
subsection. \ For now, we wish to explain how this construction relates to
prior work. In \cite{bt13}, for any constant $\delta >0$, a dual
witness is given for the fact that the $(1-\delta )$-approximate degree
of $\mathsf{OR}$ is $\Omega (\sqrt{N})$. \ This dual witness \emph{nearly}
corresponds to the above, with $w=1$. \ Specifically, Bun and Thaler \cite{bt13} use the
set $T=\{0,1\}\cup \{ci^{2}\colon i=1,2,\dots ,\sqrt{N/c}\}$, and they show
that almost all of the \textquotedblleft mass\textquotedblright\ of this
dual witness is located on the inputs $0$ and $1$, i.e.,
\begin{equation}
\left\vert \Phi (0)\right\vert +\left\vert \Phi (1)\right\vert \geq
(1-\delta )\cdot \sum_{i=2}^{N}\left\vert \Phi (i)\right\vert .
\label{epseq}
\end{equation}
Here, the bigger $c$ is chosen to be, the smaller the value of $\delta$
for which \Cref{epseq} holds.

In the case of $w=1$, our dual witness for approximate counting differs from
this only in that $\{0, 1\}$ is replaced with $\{1, 2\}$. This is
because, in order to show a lower bound for distinguishing input $w=1$ from
input $2w=2$, we want almost all of the mass to be on inputs $\{1, 2\}$
rather than $\{0, 1\}$ (this is what will ensure that the objective function
of the dual linear program is large).

For general $w$, we want most of the mass of $\psi $ to be concentrated on
inputs $w$ and $2w$. \ Accordingly, relative to the $w=1$ case, we
effectively multiply \emph{all} points in $T$ by $w$, and one can show that
this does not affect the calculation regarding concentration of mass.

\para{A dual witness for purely negative degree.} \
Now, suppose we were merely interested in showing that Laurent polynomials
of \emph{purely negative} degree require degree $\Omega (w^{1/3})$ to
approximate the approximate counting problem. \ This is equivalent to
exhibiting a solution to the dual linear program with $D_{2}=0$.  Then a
valid dual witness $\phi_2 $ for this simpler case is to also use 
\Cref{phidef}, but to set

\begin{equation}  \label{simplerT2}
T= \{w, 2w\} \cup T_1.
\end{equation}

Again, we will give intuition for why this is a valid dual solution in the
next subsection (\Cref{s:nextsubsection}). \ For now, we wish to explain how
this construction relates to prior work. \ Essentially, the $\Omega \left(
w^{1/3}\right) $-degree lower bound for Laurent polynomials with \emph{only
negative} powers was proved by Zhandry \cite[Theorem 7.3]{zhandry}. \
Translating Zhandry's theorem into our setting is not entirely trivial, and
he did not explicitly construct a solution to our dual linear program. \
However (albeit with significant effort), one can translate his argument to
our setting to show that \Cref{simplerT2} gives a valid dual
solution to prove a lower bound against Laurent polynomials with only
negative powers.

\para{Gluing them together.} The above discussion
explains that the key ideas for constructing dual solutions $\phi _{1}$, $
\phi _{2}$ witnessing degree lower bounds for Laurent polynomials of \emph{
only negative} or \emph{only positive} powers were essentially already
known, or at least can be extracted from prior work with enough effort. \ In this work, we are interested in proving lower bounds for Laurent
polynomials with both positive and negative powers. \ Our dual solution $
\Phi $ essentially just \textquotedblleft glues together\textquotedblright\
the dual solutions that can be derived from prior work. \ By this, we mean
that the set $T$ of integer points on which our $\Phi $ is nonzero is the
\emph{union} of the corresponding sets for $\phi _{1}$ and $\phi _{2}$
individually. \ Moreover, this union is nearly disjoint, as the only points
in the intersection of the two sets being unioned are $w$ and $2w$.

\para{Overview of the analysis.}
To show that we have constructed
a valid solution 
to the dual linear program (\Cref{duallp}),
we
must establish
that (a) $\Phi$ is uncorrelated
with every polynomial of degree 
at most $D_1 + D_2$
and (b) $\Phi$ is
well-correlated  with 
any function $g$ that evaluates
to $+1$ on input $w$, to $-1$ 
on input $2w$, and is bounded in $[-1, 1]$ elsewhere.
In (b), the correlation is taken
with respect to 
an appropriate weighting
of the inputs, that on input $\ell \in [N]$ places
mass proportional to $\ell^{D_1}$.

The definition of $\Phi$
as a ``gluing together'' of $\phi_1$ and $\phi_2$
turns out, in a
straightforward manner, to ensure
that $\Phi$ is uncorrelated
with polynomials of degree at 
$D_1 + D_2$. 
All that remains is to
show that $\Phi$ is well-correlated
with $g$ under the appropriate
weighting of inputs. 
This turns out to be technically demanding, but ultimately
can be understood as stemming from the fact that 
$\phi_1$ and $\phi_2$ are individually
well-correlated with $g$ (albeit, in the case of $\phi_2$, under a
\emph{different} weighting of the inputs than the weighting that is relevant for $\Phi$). 

\subsubsection{Intuition via complementary slackness}

\label{s:nextsubsection}We now attempt to lend some insight into why the
dual witnesses $\phi_1 $ and $\phi_2 $ for the purely positive degree and purely
negative degree take the form that they do. \ This section is deliberately
slightly imprecise in places, and builds on intuition that has been put
forth in prior works proving approximate degree lower bounds via dual
witnesses \cite{bt13, Tha16, bkt}.

Notice that $\phi_1 $ is precisely defined so that $\phi_1 (i)=0$ for any $
i\not\in \{w,2w\}\cup T_{2}$, and similarly $\phi_2 (i)=0$ for any $i\not\in
\{w,2w\}\cup T_{1}$. \ The intuition for why this is reasonable comes from
complementary slackness, which states that an optimal dual witness should
equal $0$ except on inputs that correspond to primal constraints that are
\emph{made tight by an optimal primal solution}. \ By \textquotedblleft
constraints made tight by an optimal primal solution\textquotedblright , we
mean constraints that, for the optimal primal solution, hold with equality
rather than (strict) inequality.

Unpacking that statement, this means the following. \ Suppose that $q$ is an
optimal solution to the primal linear program of \Cref{s:duallp}, meaning it
minimizes the error $\varepsilon $ amongst all polynomials of the same same
degree. \ The constraints made tight by $q$ are precisely those inputs $\ell
$ at which $q$ hits its \textquotedblleft maximum error\textquotedblright\
(e.g., an input $\ell$ such that $|q(\ell)|=(1+\varepsilon )\cdot \ell ^{D_{1}}$).
\ We call these inputs \emph{maximum-error} inputs for $q$. \ Complementary
slackness says that there is an optimal solution to the dual linear program (\Cref{duallp})
that equals $0$ at all inputs that are not maximum-error inputs for $q$.

In both the purely positive degree case, and the purely negative degree
case, we know roughly what primal optimal solutions $q$ look like, and
moreover we know what roughly their maximum-error points look like. \ In the
first case, the maximum-error points are well-approximated by the points in $
T_{2}$, and in the purely negative degree case, the maximum error points are
well-approximated by the points in $T_{1}$. \ Let us explain.

\para{Purely positive degree case.} \ Let $T_{d}$ be
the degree $d$ Chebyshev polynomial of the first kind.  It can be seen that $
P(\ell )=T_{\sqrt{N}}\left( 1+2/N-\ell /N\right) $ satisfies $P(1)\geq 2$,
while $|P(\ell )|\leq 1$ for $\ell =2,3,\dots ,N$. \ That is, up to scaling,
$P$ approximates the approximate counting problem for $w=1$, and its known
that its degree is within a constant factor of optimal.

It is known that the extreme points of $T_{d}$ are of the following form,
for $k=1,\dots ,d$:
\begin{equation}
\cos \left( \frac{(2k-1)}{2d}\pi \right) \approx 1-k^{2}/(2d^{2}),
\label{tester}
\end{equation}
where the approximation uses the Taylor expansion of the cosine function
around $0$. \ \Cref{tester} means that the extreme points of $P$
are roughly those inputs $\ell $ such that $1+2/N-\ell /N\approx
1-k^{2}/(2d^{2})$, where $d=\sqrt{N}$. \ Such $\ell $ are roughly of the
form $\ell \approx c\cdot i^{2}$ for some constant $c$, as $i$ ranges from $1
$ up to $\Theta (N^{1/2})$.

More generally, when $w\geq 1$, an asymptotically optimal approximation for
distinguishing input $w$ from $2w$ is $P(\ell )=T_{\sqrt{N/w}}\left(
1+2w/N-\ell /(wN)\right) $. \ The extreme points of $P$ are roughly of the
form $\ell \approx c\cdot i^{2}\cdot w$ for some constant $c$, as $i$ ranges
from $1$ up to $\Theta (\sqrt{N/w})$, which is exactly the form of the
points in our set $T_{2}$.

\para{Purely negative degree case.} In \Cref{fedjatight}
, we exhibited a simple, purely negative degree Laurent polynomial $p$
(i.e., $p(\ell )$ is a standard polynomial in $1/\ell $) with degree $
D_{1}=w^{1/3}$ that solves the approximate counting problem (the
construction is due to MathOverflow user \textquotedblleft
fedja\textquotedblright ). \ Roughly speaking, $p$ can be written as a
product $p(\ell )=u(\ell )\cdot v(\ell )$, where $u(\ell )$ has the roots $
\ell =1,2,\dots ,w^{1/3}$, and $v(\ell )$ is (an affine transformation) of a
Chebyshev polynomial of degree $w^{1/3}$, applied to $1/\ell $. \ One can
easily look at this construction and see that $p(\ell )$ outputs \emph{
exactly} the correct value on inputs $\{1,2,\dots ,w^{1/3}\}$, so these are
not maximum error points for $p$. \ Moreover, the analysis of the maximum
error points for Chebyshev polynomials above can be applied to show that the
maximum error points of $p$ are roughly of the form $\ell $ such that $
1/\ell =c\cdot i^{2}/w$ for some constant $c$, with $i$ ranging from $1$ up
to $\Theta (w^{1/3})$. \ This means that the extreme points are roughly of
the form $\ell \approx \frac{w}{ci^{2}}$, which is why our set $T_{1}$
consists of points of the form $\lfloor \frac{w}{ci^{2}}\rfloor $ (the
floors are required because we are proving lower bounds against polynomials
whose behavior is only constrained at integer inputs).

\subsubsection{Analysis of the dual solution \texorpdfstring{$\Phi$}{Phi}}

\begin{lemma}
\label{phdlem}Let $d_{1}=|T_{1}|$ and $d_{2}=|T_{2}|$. \ Then for any $
j=0,1,\dots ,d_{1}+d_{2}$, it holds that
\begin{equation*}
\sum_{\ell =1}^{N}\Phi (\ell )\cdot \ell ^{j}=0.
\end{equation*}
\end{lemma}

\begin{proof}
A basic combinatorial fact is that for any polynomial $Q$ of degree at most $N-1$,
the following identity holds:
\begin{equation} 
\label{keyidentity} 
\sum_{\ell=0}^N \binom{N}{\ell} (-1)^\ell Q(\ell) = 0.
\end{equation}

Observe that for any $j \leq d_1 + d_2+1$,
\begin{equation} \label{isapolynomial}  
Q_T(\ell) \cdot \ell^j  \text{ is a polynomial in } \ell \text{ of degree at most } N-1. 
\end{equation}
Furthermore, $\Phi(0)=0$, because $0 \not\in T$. Hence
\begin{equation} \label{nozeroterm} 
\sum_{\ell=0}^N \binom{N}{\ell} (-1)^\ell Q_T(\ell) \cdot \ell^j= \sum_{\ell=1}^N \binom{N}{\ell} (-1)^\ell Q_T(\ell) \cdot \ell^j. 
\end{equation}
Thus, we can calculate:
\begin{eqnarray*}
\sum_{\ell=1}^N \Phi(\ell) \cdot \ell^{j} = \sum_{\ell=1}^N (-1)^{\ell} \cdot \binom{N}{\ell} \cdot Q_T(\ell) \cdot \ell^j\\
= \sum_{\ell=0}^N (-1)^\ell \cdot \binom{N}{\ell} \cdot Q_T(\ell) \cdot \ell^j
=0.
\end{eqnarray*}
Here, the second equality follows from \Cref{nozeroterm}, while the third follows
from Equations \eqref{keyidentity} and \eqref{isapolynomial}.
\end{proof}

Let us turn to analyzing $\Phi $'s value on various inputs. \ Clearly the
following condition holds:

\begin{equation}
\Phi (\ell )=0~\text{for all }\ell \not\in T\text{.}  \label{zeros}
\end{equation}
Next, observe that for any $r\in T$,
\begin{equation*}
\left\vert \Phi (r)\right\vert =N!\cdot \frac{1}{\prod_{j\in T,j\neq r}|r-j|}
.
\end{equation*}

Consider any quantity $c\cdot i^{2}\cdot w\in T_{2}$. \ Then
\begin{flalign} & \left|\Phi(c \cdot w \cdot i^2)\right|/\left|\Phi(w)\right| = \frac{\prod_{j \in T, j \neq w} |w- j|}{ \prod_{j \in T, j \neq c \cdot i^2 \cdot w} |w \cdot c \cdot i^2 - j|} \notag  \\
& = \frac{\left|w-2w\right| \cdot \left(\prod_{j=1}^{d_2} \left|w-c \cdot j^2 \cdot w\right|\right) \cdot \left(\prod_{j=1}^{d_1} \left(w-\left\lfloor \frac{w}{cj^2}\right\rfloor\right)\right)}{\left|c \cdot i^2 \cdot w - w\right| \cdot \left|c \cdot i^2 \cdot w - 2w\right| \cdot \left( \prod_{j=1, j\neq i}^{d_2} \left|w \cdot c \cdot i^2 - w \cdot c \cdot j^2\right| \right)\cdot \left( \prod_{j=1}^{d_1} \left(w \cdot c \cdot i^2 - \left\lfloor\frac{w}{c \cdot j^2}\right\rfloor\right)\right)}\notag \\
& = \frac{c^{d_2} \cdot \left(\prod_{j=1}^{d_2} \left( j^2 - \frac{1}{c}\right)\right) \cdot \prod_{j=1}^{d_1} \left(w-\left\lfloor\frac{w}{c\cdot j^2}\right\rfloor\right)}{ \left(ci^2 -1\right)\cdot \left(ci^2 -2 \right) \cdot c^{d_2-1} \cdot \left( \prod_{j=1, j \neq i}^{d_2} \left|i^2 - j^2\right|\right) \cdot \left(\prod_{j=1}^{d_1} \left(w \cdot c \cdot i^2 - \left\lfloor \frac{w}{c \cdot j^2}\right\rfloor\right)\right)} \notag \\
& \leq
\frac{ c \cdot \left( \prod_{j=1}^{d_2} \left( j^2 - \frac{1}{c}\right)\right) \cdot \prod_{j=1}^{d_1} \left(w-\left\lfloor\frac{w}{c\cdot j^2}\right\rfloor\right)}{ \left(ci^2 -1\right)\cdot \left(ci^2 -2 \right) \cdot \left(\prod_{j=1, j \neq i}^{d_2} \left|i^2 - j^2\right|\right) \cdot \left(\prod_{j=1}^{d_1} \left(w \cdot c \cdot i^2 -  \frac{w}{c \cdot j^2}\right)\right)} \label{firstlongequation}
\end{flalign}

Now, observe that
\begin{align}
\prod_{j=1}^{d_1} \left(w-\left\lfloor\frac{w}{c\cdot j^2}\right\rfloor\right)
\leq \prod_{j=1}^{d_1} \left(w - \frac{w}{cj^2} +1 \right) =
\prod_{j=1}^{d_1} w \cdot \left(1-\frac{1}{cj^2}\right) \cdot \left( 1+
\frac{1}{w \cdot \left(1-\frac{1}{cj^2}\right)}\right)  \notag \\
\leq \prod_{j=1}^{d_1} w \cdot \left(1-\frac{1}{cj^2}\right) \left(1 + \frac{
1}{(1-1/c) \cdot w}\right) \leq \left( \prod_{j=1}^{d_1} w \cdot \left(1-
\frac{1}{cj^2}\right)\right) \cdot \left(1+o(1) \right).  \label{alongequation}
\end{align}

Hence, we see that Expression \eqref{firstlongequation} is bounded by

\begin{align}
& \frac{c\cdot \left( \prod_{j=1}^{d_{2}}\left( j^{2}-\frac{1}{c}\right)
\right) \cdot \left( \prod_{j=1}^{d_{1}}\left( 1-\frac{1}{c\cdot j^{2}}
\right) \right) \cdot (1+o(1))}{\left( ci^{2}-1\right) \cdot \left(
ci^{2}-2\right) \cdot \left( \prod_{j=1,j\neq i}^{d_{2}}\left\vert
i^{2}-j^{2}\right\vert \right) \cdot \left( \prod_{j=1}^{d_{1}}\left( c\cdot
i^{2}-\frac{1}{c\cdot j^{2}}\right) \right) }  \notag \\
& \leq \frac{c\cdot \left( d_{2}!\right) ^{2}\cdot \left(
\prod_{j=1}^{d_{1}}\left( 1-\frac{1}{c\cdot j^{2}}\right) \right) \cdot
(1+o(1))}{\left( ci^{2}-1\right) \cdot \left( ci^{2}-2\right) \cdot \left(
\prod_{j=1,j\neq i}^{d_{2}}\left\vert i-j\right\vert \left\vert
i+j\right\vert \right) \cdot (c\cdot i^{2})^{d_{1}}\cdot \left(
\prod_{j=1}^{d_{1}}\left( 1-\frac{1}{c^{2}\cdot i^{2}\cdot j^{2}}\right)
\right) }  \notag \\
& =\frac{c\cdot \left( d_{2}!\right) ^{2}\cdot 2i^{2}\cdot \left(
\prod_{j=1}^{d_{1}}\left( 1-\frac{1}{c\cdot j^{2}}\right) \right) \cdot
(1+o(1))}{\left( ci^{2}-1\right) \cdot \left( ci^{2}-2\right) \cdot \left(
d_{2}+i\right) !\left( d_{2}-i\right) !\cdot (c\cdot i^{2})^{d_{1}}\cdot
\left( \prod_{j=1}^{d_{1}}\left( 1-\frac{1}{c^{2}\cdot i^{2}\cdot j^{2}}
\right) \right) }  \notag \\
& \leq \frac{c\cdot 2i^{2}\cdot \left( d_{2}!\right) ^{2}\cdot (1+o(1))}{
\left( ci^{2}-1\right) \left( ci^{2}-2\right) \cdot \left( d_{2}+i\right)
!\left( d_{2}-i\right) !\cdot (c\cdot i^{2})^{d_{1}}}\leq \frac{2\left(
1+o(1)\right) }{\left( 1-\frac{1}{c\cdot i^{2}}\right) \cdot (c\cdot
i^{2}-2)\cdot (c\cdot i^{2})^{d_{1}}}.  \label{verylongequation}
\end{align}
In the penultimate inequality, we used the fact that $\frac{(d_{2}!)^{2}}{
(d_{2}+i)!(d_{2}-i)!}=\frac{\binom{2d_{2}}{d_{2}+i}}{\binom{2d_{2}}{d_{2}}}
\leq 1$.

Next, consider any quantity $\left\lfloor \frac{w}{c\cdot i^{2}}\right\rfloor \in T_{1}$
. \ Then
\begin{flalign} & \left|\Phi\left( \left\lfloor \frac{w}{c \cdot i^2}\right\rfloor\right)\right|/\left|\Phi(w)\right| \notag \\
& = \frac{|w-2w| \left(\prod_{j=1}^{d_2} |w - cj^2 w| \right) \left(\prod_{j=1}^{d_1} \left( w - \left\lfloor  \frac{w}{cj^2}\right\rfloor\right)\right)}{ \left(w-\left\lfloor \frac{w}{c \cdot i^2}\right\rfloor\right) \cdot \left(2w - \left\lfloor \frac{w}{c \cdot i^2}\right\rfloor \right) \left(\prod_{j=1}^{d_2} \left(w \cdot c \cdot j^2 - \left\lfloor \frac{w}{c \cdot i^2}\right\rfloor\right)\right) \prod_{j=1, j\neq i}^{d_1} \left|\left\lfloor \frac{w}{c \cdot i^2}\right\rfloor - \left\lfloor \frac{w}{c \cdot j^2}\right\rfloor\right|} \notag \\
& \leq
\frac{|w-2w| \left(\prod_{j=1}^{d_2} |w - cj^2 w| \right) \left(\prod_{j=1}^{d_1} \left( w - \left\lfloor  \frac{w}{cj^2}\right\rfloor\right)\right)}{ \left(w- \frac{w}{c \cdot i^2}\right) \cdot \left(2w -  \frac{w}{c \cdot i^2} \right) \left(\prod_{j=1}^{d_2} \left(w \cdot c \cdot j^2 - \frac{w}{c \cdot i^2}\right)\right) \prod_{j=1, j\neq i}^{d_1} \left|\left\lfloor \frac{w}{c \cdot i^2}\right\rfloor - \left\lfloor \frac{w}{c \cdot j^2}\right\rfloor\right|} \notag \\
& \leq  \frac{|w-2w| \left(\prod_{j=1}^{d_2} |w - cj^2 w| \right) \left(\prod_{j=1}^{d_1} \left( w -  \frac{w}{cj^2}\right)\right) \cdot \left(1+o(1)\right)}{ \left(w- \frac{w}{c \cdot i^2}\right) \cdot \left(2w -  \frac{w}{c \cdot i^2} \right) \left(\prod_{j=1}^{d_2} \left(w \cdot c \cdot j^2 - \frac{w}{c \cdot i^2}\right)\right) \prod_{j=1, j\neq i}^{d_1} \left|\left\lfloor \frac{w}{c \cdot i^2}\right\rfloor - \left\lfloor \frac{w}{c \cdot j^2}\right\rfloor\right|} \label{anotherlongequation}
\end{flalign}
Here, the final inequality used \Cref{alongequation}. 
Let us consider the expression 
$\prod_{j=1,j\neq i}^{d_{1}}\left\vert \left\lfloor \frac{
w}{c\cdot i^{2}}\right\rfloor - \left\lfloor \frac{w}{c\cdot j^{2}}\right\rfloor \right\vert $
. \ 
This quantity is \emph{at least}
\begin{align}
\prod_{j=1, j\neq i}^{d_1}\left( \left|\frac{w}{c \cdot i^2} - \frac{w}{c
\cdot j^2} \right| -1\right) = w^{d_1-1} \cdot \prod_{j=1, j \neq i}^{d_1}
\frac{\left|j^2 - i^2\right| - \frac{c i^2 j^2}{w}}{ci^2 j^2}  \notag \\
= w^{d_1-1} \cdot \prod_{j=1, j \neq i}^{d_1} \frac{\left|j-i| \cdot
|j+i\right| - \frac{c i^2 j^2}{w}}{ci^2 j^2}  \notag \\
= \left(\frac{w}{ci^2}\right)^{d_1-1} \cdot \prod_{j=1, j \neq i}^{d_1}
\frac{\left|j-i| \cdot |j+i\right| - \frac{c i^2 j^2}{w}}{j^2}
\label{aboveexp}
\end{align}

We claim that Expression \eqref{aboveexp} is at least
\begin{equation}
\left( \frac{w}{ci^{2}}\right) ^{d_{1}-1}\cdot \frac{1}{2}.
\label{provedintheappendix}
\end{equation}
In the case that $c=2$ and $d_{1}$ is (at most) $w^{1/3}$, this is precisely
\cite[Claim 4]{zhandry}. \ We will ultimately take $c$ to be a constant
strictly greater than 2 and hence $d_{1}=\lfloor \left( w/c\right)
^{1/3}\rfloor $ is a constant factor smaller than $w^{1/3}$. \ The proof of
\cite[Claim 4]{zhandry} works with cosmetic changes in this case. \ For
completeness, we present a derivation of the claim in \Cref{app}.

\Cref{provedintheappendix} implies that Expression \eqref{anotherlongequation} is
at most:

\begin{align}
\frac{|w-2w| \left(\prod_{j=1}^{d_2} |w - cj^2 w| \right)
\left(\prod_{j=1}^{d_1} \left( w - \frac{w}{cj^2}\right)\right) \cdot
\left(1+o(1)\right)}{ \left(w- \frac{w}{c \cdot i^2}\right) \cdot \left(2w -
\frac{w}{c \cdot i^2} \right) \left(\prod_{j=1}^{d_2} \left(w \cdot c \cdot
j^2 - \frac{w}{c \cdot i^2}\right)\right) \left(\frac{w}{ci^2}
\right)^{d_1-1} \cdot \frac{1}{2}}  \notag \\
= \frac{2 \left(\prod_{j=1}^{d_2} |1 - cj^2| \right) \left(\prod_{j=1}^{d_1}
\left(1 - \frac{1}{cj^2}\right)\right) \cdot \left(1+o(1)\right)}{ \left(1-
\frac{1}{c \cdot i^2}\right) \cdot \left(2 - \frac{1}{c \cdot i^2} \right)
\left(\prod_{j=1}^{d_2} \left(c \cdot j^2 - \frac{1}{c \cdot i^2}
\right)\right) \left(\frac{1}{ci^2}\right)^{d_1-1}}  \notag \\
= \frac{2 \left(\prod_{j=1}^{d_2} (j^2-1/c) \right) \left(\prod_{j=1}^{d_1}
\left(1 - \frac{1}{cj^2}\right)\right) \cdot \left(1+o(1)\right)}{ \left(1-
\frac{1}{c \cdot i^2}\right) \cdot \left(2 - \frac{1}{c \cdot i^2} \right)
\left(\prod_{j=1}^{d_2} \left(j^2 - \frac{1}{c^2 \cdot i^2}\right)\right)
\left(\frac{1}{ci^2}\right)^{d_1-1}}  \notag \\
\leq \frac{2 \left(1+o(1)\right)}{ \left(1- \frac{1}{c \cdot i^2}\right)
\cdot \left(2 - \frac{1}{c \cdot i^2} \right) \left(\frac{1}{ci^2}
\right)^{d_1-1}} \leq 4 \cdot \left(ci^2\right)^{d_1-1}.  \label{implication}
\end{align}

Summarizing Equations \eqref{verylongequation} and \eqref{implication}, we have shown
that: for any quantity $c \cdot i^2 \cdot w \in T_2$,
\begin{equation}  \label{keyeq1}
\left|\Phi(c \cdot w \cdot i^2)\right|/\left|\Phi(w)\right| \leq \frac{2
\left(1+o(1)\right)}{\left(1-\frac{1}{c \cdot i^2}\right) \cdot (c \cdot i^2
- 2) \cdot (c\cdot i^2)^{d_1}}
\end{equation}
and for any quantity $\left\lfloor \frac{w}{c \cdot i^2} \right\rfloor \in T_1$,
\begin{equation}  \label{keyeq2}
\left|\Phi\left( \left\lfloor \frac{w}{c \cdot i^2}\right\rfloor\right)\right|/\left|\Phi(w)
\right| \leq 4\cdot \left(ci^2\right)^{d_1-1}.
\end{equation}

Let $\phi =\Phi /C$, where $C$ is as in \Cref{Cdef}. \ Let $
D_{1}=d_{1}$ and $D_{2}=d_{2}$. \ \Cref{phdlem} implies that $\phi $ is a
feasible solution for the dual linear program of \Cref{duallpsec}. \ We now
show that, for any constant $\delta >0$, by choosing $c$ to be a
sufficiently large constant (that depends on $\delta $), we can ensure
that $\phi $ achieves objective value $1-2\delta $.

Let
\begin{eqnarray*}
A &=&|\Phi (w)|\cdot w^{D_{1}}, \\
B &=&|\Phi (2w)|\cdot (2w)^{D_{1}},
\end{eqnarray*}
and
\begin{equation*}
E=\sum_{i=1}^{d_{1}}|\Phi (\lfloor w/ci^{2}\rfloor )|\cdot \left( \lfloor
w/ci^{2}\rfloor \right) ^{D_{1}}+\sum_{i=1}^{d_{2}}|\Phi (\lfloor w\cdot
ci^{2}\rfloor )|\cdot \left( w\cdot c\cdot i^{2}\right) ^{D_{1}}.
\end{equation*}
By \Cref{zeros}, $C=A+B+E$.

Moreover, observe that $\text{sgn}(\Phi(w)) = - \text{sgn}(\Phi(2w))$, so
without loss of generality we may assume $\Phi(w) \geq 0$ and $\Phi(2w) \leq
0$ (if not, then replace $\Phi$ with $-\Phi$ throughout).

We now claim that by choosing $c$ to be a sufficiently large constant, we
can ensure that $E \leq \delta \cdot A$. To see this, observe that
Equations \eqref{keyeq1} and \eqref{keyeq2}, along with the fact that $
D_1=d_1$ and $D_2=d_2$ implies that
\begin{align*}
E/A \leq \frac{1}{w^{D_1}} \left[ \left(\sum_{i=1}^{d_1} \left(\lfloor
w/ci^2\rfloor\right)^{D_1}\cdot 4 \cdot \left(ci^2\right)^{d_1-1} \right) +
\left(\sum_{i=1}^{d_2} \left(w \cdot c \cdot i^2\right)^{D_1} \frac{2
\left(1-\frac{1}{c \cdot i^2}\right)\left(1+o(1)\right)}{(c \cdot i^2 - 2)
\cdot (c\cdot i^2)^{d_1}}\right)\right] \\
\leq \frac{1}{w^{D_1}} \left[ \left(\sum_{i=1}^{d_1}
\left(w/ci^2\right)^{D_1} \cdot 4 \cdot \left(ci^2\right)^{d_1-1} \right) +
\left(\sum_{i=1}^{d_2} \left(w \cdot c \cdot i^2\right)^{D_1} \frac{2
\left(1-\frac{1}{c \cdot i^2}\right)\left(1+o(1)\right)}{(c \cdot i^2 - 2)
\cdot (c\cdot i^2)^{d_1}}\right) \right] \\
\leq 4 \left(\sum_{i=1}^{d_1} \frac{1}{c\cdot i^2}\right) +
\left(\sum_{i=1}^{d_2} \frac{2 \left(1+o(1)\right)}{\left(1-\frac{1}{c \cdot
i^2}\right) \left(c \cdot i^2 - 2\right)}\right) \\
\end{align*}

Since $\sum_{i=1}^{\infty}1/(ci^2) \leq \frac{\pi^2}{6c}$, we see that
choosing $c$ to be a sufficiently large constant depending on $\delta$
ensures that $E/A \leq \delta$ as desired.

Hence, $\phi$ achieves objective value at least
\begin{align*}
\phi(w) \cdot w^{D_1} - \phi(2w) \cdot (2w)^{D_1} - \sum_{\ell \in \{1,
\dots, N\}, \ell \not\in \{w, 2w\}} |\phi(\ell)| \cdot \ell^{D_1} \\
\geq \frac{A + B - E}{A+B+E} \geq \frac{(1-\delta) A + B}{
(1+\delta)A + B} \geq 1-2\delta.
\end{align*}

%
\setcounter{subsection}{5}
\setcounter{theorem}{30}
\subsection{Extending the lower bound to QSampling unitarily}
\label{sec:unitary}
So far in this section we have proved upper and lower bounds on the power of quantum algorithms for approximate counting that have 
access to two resources (in addition to membership queries): 
copies of $|S\>$, and the unitary transformation that reflects about $|S\>$. 
The assumption of access to the reflection unitary is justified by the argument that, if we had access to a unitary that prepared $|S\>$, then it could be used to reflect about $|S\>$ as well. 

Giving the algorithm access 
to just the two resources above is an appealing model to use for upper bounds, since it does not assume anything about the method by which copies of $|S\>$ are prepared. This means algorithms
derived in this model work in many different settings.
For example, the algorithm may be able to QSample
because someone else simply handed the algorithm copies of $|S\>$, or perhaps several copies of $|S\>$ just happen to be stored in the algorithm's quantum memory as a side effect of the execution of some earlier quantum algorithm. The upper bound given
in \Cref{thm:alg} applies in any of these settings.


On the other hand, since only
permitting access to QSamples
and reflections about $|S\>$ ties the algorithm's hands,
lower bounds for this model
(e.g., \Cref{thm:main}) could be viewed
as weaker than is desirable. 
In particular, our original 
justification
for allowing access to reflections about $|S\>$
was that access to a unitary that prepared 
the state $|S\>$ would in particular allow 
such reflections to be done. 
Given this justification, it is very natural to wonder whether 
our lower bounds extend beyond just 
QSamples and reflections, to
algorithms that are given access
to \emph{some} unitary process that 
permits both QSampling and reflections about $|S\>$. 

Note that
an algorithm with access to such a unitary could potentially exploit the unitary in ways other
than QSamples and reflections to learn information about $|S\>$. 
For example, the algorithm could choose to run the unitary on inputs that do not produce $|S\>$.
More generally, given a quantum circuit that implements a unitary, it is possible to construct, in a completely black-box manner, the inverse of this unitary, and also a controlled version of the unitary.
The algorithm may choose to run the inverse on a state other than $|S\>$ to learn some additional information that is not captured by access
to QSamples and reflections alone.

In summary, in this section
we ask whether we can we extend the lower bound of \Cref{thm:main} to work in a model where the algorithm is given access to some unitary operator that conveys the power to both QSample and reflect about $|S\>$.\footnote{We thank Alexander Belov (personal communication) for raising this question.} 
Via \cref{thm:finalthm} below, we explain that the answer is yes.


\medskip

It may seem convenient to assume that the unitary
transformation preparing $|S\>$ 
maps the all-zeros state to $|S\>$.
But this is not the most general method of preparing $|S\>$ by a unitary.
A unitary $U$ that maps the all-zeros state to $|S\>|\psi\>$ would also suffice to create copies of $|S\>$, since the register containing $|\psi\>$ can simply be ignored for the remainder of the computation.
More formally, assume $U$ behaves as 
\begin{equation}
    U|0^m\> = |S\>|\psi\>,
\end{equation}
where $|S\>|\psi\>$ is some $m$-qubit state.
Clearly we can use $U$ to create as many copies of $|S\>$ as we like, 
which as a by-product also creates copies of $|\psi\>$.
This unitary also lets us reflect about $|S\>$. 
To see how, first use this unitary to create a copy of $|\psi\>$, and then consider the action of the unitary $U(\id - 2|0^m\>\<0^m|)U^\dagger$ on the state $|\phi\>|\psi\>$ for any state $|\phi\>$. We claim that this unitary acts as a reflection about $|S\>$ when restricted to the first register. 
This establishes that any $U$ of this form subsumes the power of both QSamples and reflections about $|S\>$.

Let us also assume without loss of generality that $|S\>|\psi\>$ is orthogonal to $|0^m\>$ from now on. This can be achieved by adding an additional qubit to the input that is always negated by the unitary. That is, we could instead consider the map $(U \otimes X)|0^m\>|0\> = |S\>|\psi\>|1\>$, which is orthogonal to the starting state by construction, and only increases the value of $m$ by $1$.

Of course, the requirement that $U|0^m\> = |S\>|\psi\>$ does not fully specify $U$, as
it does not prescribe how $U$ behaves on other input states.
A reasonable prescription is that $U$ should behave ``trivially'' on other input states, so that it does not leak information about $S$ by its behavior on other states. 
In tension with this prescription is the fact the rest of the unitary must depend on $S$, since the first column of the unitary contains $|S\>$, and the rest of the columns have to be orthogonal to this.

Alexander Belov (personal communication) brought to our attention a 
very simple construction of such a unitary that leaks minimal additional information about $S$. 
Consider the unitary $U$ that satisfies $U|0^m\> = |S\>|\psi\>$ and $U|S\>|\psi\> = |0^m\>$, with $U$ acting as identity outside $\mathrm{span}\{|0^m\>,|S\>|\psi\>\}$. $U$ is simply a reflection about the state $\frac{1}{\sqrt{2}}\bigl(|0^m\>-|S\>|\psi\>\bigr)$. This state is correctly normalized because we assumed that $|S\>|\psi\>$ is orthogonal to $|0^m\>$. 
Clearly $U$ is now fully specified on the entire domain (once we have fixed $|\psi\>$) and it does not seem to leak any additional information about $S$.

In order to prove concrete lower bounds on the cost
of algorithms for approximate counting given access to $U$, we need to fix $|\psi\>$. 
To answer the question posed in this section, we only need to establish that there exists \emph{some} choice of $|\psi\>$ for which our algorithms cannot be improved. (Note that we cannot hope to establish lower bounds for arbitrary $|\psi\>$, since $|\psi\>$ could just contain the answer to the problem we are solving.)

To this end we make the specific choice of $|\psi\> = |S\>$ and consider the unitary $V$ that acts as the unitary $U$ above with $|\psi\> = |S\>$. 
In other words, $V$ maps $|0^m\>$ to $|S\>|S\>$, $|S\>|S\>$ to $|0^m\>$, and acts as identity on the rest of the space. 
We also assume that $|0^m\>$ is orthogonal to $|S\>|S\>$. 
In other words, $V$ simply reflects about the state $\frac{1}{\sqrt{2}}\bigl(|0^m\>-|S\>|S\>\bigr)$.

As previously discussed, granting an algorithm access to this unitary $V$ lends the algorithm at least as much power the ability to QSample and perform reflections about $|S\>$. How efficiently can we solve approximate counting with membership queries and uses of the unitary $V$? 

We can use our Laurent polynomial method to establish optimal lower bounds in this model as well and we obtain lower bounds identical to \Cref{thm:main}. 

\begin{theorem} \label{thm:finalthm} \label{thm:unitary}
Let $Q$ be a quantum algorithm that makes $T$ queries to the membership oracle for $S$,
and makes $R$ uses of the unitary $V$ defined above (and its inverse and controlled-$V$).
If $Q$ decides whether $\left\vert S\right\vert =w$\ or $\left\vert S\right\vert =2w$
with high probability, promised that one of those is the case, then either
\begin{equation}
T=\Omega\left(  \sqrt{\frac{N}{w}}\right)  \qquad \textrm{or} \qquad
R=\Omega\left(  \min\left\{  w^{1/3},\sqrt{\frac{N}{w}}\right\}  \right).
\end{equation}
\end{theorem}
\begin{proof}
We follow the same strategy as in
the proof of \Cref{thm:main}. Recall that $x \in \{0, 1\}^N$ denotes the indicator vector of the set $S$. We only need to show that such a quantum algorithm gives rise to a Laurent polynomial in $|S|:=\sum_{i=1}^n x_i$, with maximum exponent $O(T+R)$ and minimum exponent at least $-O(R)$ (as shown in \Cref{laurentlem} for the QSamples and reflections model).

We can prove this exactly the same way as \Cref{laurentlem} is established. Our quantum algorithm starts out from a canonical starting state that does not depend on the input and hence each entry of the starting state is a degree-$0$ polynomial.
Membership queries involve multiplication with an oracle whose entries are ordinary polynomials of degree at most $1$. The only thing that remains is understanding what the entries of the unitary $V$ look like. We claim that the entries of $V$ are given
by a polynomial of degree at most 2 in the entries of the input $x$,
with all coefficients of this degree-2 polynomial equal
to either a constant, or a constant multiple of $|S|^{-1}$.

To see this, note that $V$ is simply a reflection about the state
\begin{equation}
    \frac{1}{\sqrt{2}}\bigl(|0^m\>-|S\>|S\>\bigr)
    = \frac{1}{\sqrt{2}}\left(|0^m\>-\frac{1}{|S|}\Bigl(\sum_{i}x_i|i\>\Bigr)\Bigl(\sum_{j}x_j|j\>\Bigr)
    \right).
\end{equation} 
The coefficient in front of $|0^m\>$ is a degree-$0$ polynomial and the other nonzero coefficients are
a polynomial of degree at most 2 in the entries of the input $x$, with each coefficient of this polynomial equal to a constant multiple of $|S|^{-1}$.

Hence, each entry of the unitary $V$ is also a polynomial of degree at most 2 in the entries of the input $x$,
with each coefficient of this degree-2 polynomial equal
to either a constant, or a constant multiple of $|S|^{-1}$. The same also holds for controlled-$V$, since that unitary is just the direct sum of identity with $V$. $V$ is also self-inverse, so we do not need to account for that separately.

After the algorithm has made all the membership queries and uses of $V$, each amplitude of the final quantum state can be expressed as a polynomial of degree $O(T + R)$ in the input $x$, in which all coefficients are constant multiples of $|S|^{-R}$. The acceptance probability $p(x)$ of this algorithm will be a sum of squares of such polynomials. 
Exactly
as in the proof of \Cref{thm:main},
\Cref{symlem} implies that there is a univariate polynomial
$q$ of degree at most $O(T+R)$, with coefficients that are multiples
of the coefficients of $p$, such that for all integers $k \in \{0, \dots, N\}$, 
\begin{equation}
q\left(  k\right)  :=\E_{\left\vert X\right\vert =k}\left[
p\left(  X\right)  \right]  .
\end{equation}
Since the coefficients of $p(X)$ are constant multiples of $|X|^{-2R}$, 
$q$ is in fact a real Laurent polynomial in $k$, with maximum exponent at most $O(R+T)$\ and minimum exponent at
least $-2R$.
The theorem follows by a direct 
application \Cref{mainthm} to $q$.
\end{proof}
%

\appendix
\section{Establishing Equation \protect\ref*{provedintheappendix}}

\label{app}

%
\subsection{A clean calculation establishing a loose version of equation
\protect\ref*{provedintheappendix}}

For clarity of exposition, we begin by presenting a relatively clean
calculation that establishes a slightly loose version of 
\Cref{provedintheappendix}. Using just this looser bound, we would be able to
establish that \Cref{provedintheappendix} holds (with the constant $1/2$
replaced by a slightly smaller constant) so long as we set $d_1$ to be $
\Theta\left(w^{1/3}/\log w \right)$. A slightly more involved calculation
(cf. \Cref{tightsec}) is required to establish 
\Cref{provedintheappendix} for our desired value of $d_1= \lfloor (w/c)^{1/3}
\rfloor$.

Expression \eqref{aboveexp} equals
\begin{align}
\left(\frac{w}{ci^2}\right)^{d_1-1} \cdot \frac{i^2}{\left((d_1)!\right)^2}
\cdot \prod_{j=1, j \neq i}^{d_1}\left( \left|j-i| \cdot |j+i\right| - \frac{
c i^2 j^2}{w}\right)  \notag \\
= \left(\frac{w}{ci^2}\right)^{d_1-1} \cdot \frac{i^2}{\left((d_1)!\right)^2}
\cdot \prod_{j=1, j \neq i}^{d_1}\left( \left|j-i| \cdot |j+i\right| \right)
\cdot \left(1-\frac{c i^2 j^2}{w \cdot |j-i| |j+i|}\right)  \notag \\
= \left(\frac{w}{ci^2}\right)^{d_1-1} \cdot \frac{(d_1+i)! (d_1-i)!}{
2\left((d_1)!\right)^2} \cdot \prod_{j=1, j \neq i}^{d_1} \left(1-\frac{c
i^2 j^2}{w \cdot |j-i| |j+i|}\right)  \label{looseversion} \\
\geq \left(\frac{w}{ci^2}\right)^{d_1-1} \cdot \frac{1}{2} \cdot \prod_{j=1,
j \neq i}^{d_1} \left(1-\frac{c i^2 j^2}{w \cdot |j-i| |j+i|}\right)  \notag
\\
\geq \left(\frac{w}{ci^2}\right)^{d_1-1} \cdot \frac{1}{2} \cdot
\left(1-\sum_{j=1, j \neq i }^{d_1}\frac{c i^2 j^2}{w \cdot |j-i| |j+i|}
\right)  \notag \\
\geq \left(\frac{w}{ci^2}\right)^{d_1-1} \cdot \frac{1}{2} \cdot \left(1-
\frac{c i^2}{w}\sum_{j=1, j \neq i }^{d_1}\frac{j^2}{ |j-i| |j+i|}\right).
\label{loosemiddle}
\end{align}

Let us consider the expression $\sum_{j=1, j \neq i }^{d_1}\frac{j^2}{ |j-i|
|j+i|}$. If $i^2 \not \in [j^2/2, 3j^2/2]$, then the $j$'th term in this sum
is at most $2$. Hence, letting $H_i$ denote the $i$th Harmonic number and
using the fact that $H_i \leq \ln(i+1)$,
\begin{align}
& \sum_{j=1, j \neq i }^{d_1}\frac{j^2}{ |j-i| |j+i|}  \notag \\
& \leq 2 \cdot d_1 + \sum_{j = \lfloor \sqrt{2/3} i \rfloor}^{\lfloor \sqrt{2
} i \rfloor} \frac{j^2}{|j-i| |j+i|}  \notag \\
& \leq 2 \cdot d_1 + \sum_{j = \lfloor \sqrt{2/3} \cdot i \rfloor}^{\lceil
\sqrt{2} \cdot i \rceil} \frac{j}{|j-i|}  \notag \\
& \leq 2 d_1 + \sqrt{2} \cdot i \cdot \sum_{j=1}^{(\sqrt{2}-1) \cdot i} 2/j
\notag \\
& \leq 2 d_1 + 2 \sqrt{2} \cdot i \cdot H_{i} \leq 2d_1 + 2 \sqrt{2} i
\ln(i+1).  \label{eq:ieq}
\end{align}

We conclude that if $d_1$ were set to a value less than $w^{1/3}/(100 \cdot
c^2 \cdot \ln(w))$ (rather than to $\lfloor \left(w/c\right)^{1/3}\rfloor$),
then Expression \eqref{loosemiddle} is at least
\begin{equation}
\left(\frac{w}{ci^2}\right)^{d_1-1} \cdot \frac{1-1/c}{2}.
\end{equation}

%
\subsection{The tight bound}

\label{tightsec} To obtain the tight bound, we need a tighter sequence of
inequalities following Expression \eqref{looseversion}. Specifically,
Expression \eqref{looseversion} is bounded below by:

\begin{align}
\geq \left(\frac{w}{ci^2}\right)^{d_1-1} \cdot \frac{1}{2} \left(1+ \frac{i}{
2d_1}\right)^i \cdot \prod_{j=1, j \neq i}^{d_1} \left(1-\frac{c i^2 j^2}{w
\cdot |j-i| |j+i|}\right)  \notag \\
\geq \left(\frac{w}{ci^2}\right)^{d_1-1} \cdot \frac{1}{2} \cdot
e^{i^2/(2d_1)} \cdot \prod_{j=1, j \neq i}^{d_1} \left(1-\frac{c i^2 j^2}{w
\cdot |j-i| |j+i|}\right)  \notag \\
\geq \left(\frac{w}{ci^2}\right)^{d_1-1} \cdot \frac{1}{2} \cdot
e^{i^2/(2d_1)} \cdot \prod_{j=1, j \neq i}^{d_1} \left(1-\frac{c i^2 j^2}{w
\cdot |j-i| |j+i|}\right)  \label{tightbegin}
\end{align}

The rough idea of how to proceed is as follows. \Cref{eq:ieq}
implies that for $i \ll w^{1/3}/\ln w$, the factor
\begin{equation*}
F_1 := \prod_{j=1, j \neq i}^{d_1} \left(1-\frac{c i^2 j^2}{w \cdot |j-i|
|j+i|}\right)
\end{equation*}
is at some a positive constant, and hence Expression \eqref{tightbegin} is
bounded below by the desired quantity. If $i \gtrsim w^{1/3}/\ln w$, then
\Cref{eq:ieq} does not yield a good bound on this factor, leaving
open the possibility that this factor is subconstant. But in this case, the
factor $F_2:= e^{i^2/(2d_1)} \geq e^{\tilde{\Omega}(d_1)}$, and the
largeness of $F_2$ dominates the smallness of $F_1$.

In more detail, let $x_{i, j} = \frac{ci^2j^2}{w \cdot |i-j| j+i||}$. Then
for all $i \neq j$ such that $i, j \leq d_1$,
\begin{equation}  \label{xijbound}
x_{i,j} \leq \frac{c \cdot d_1^2 (d_1-1)^2}{(2d_1-1) \cdot w} \leq \frac{c
\cdot d_1^3}{2w} \leq 1/2,
\end{equation}
where in the final inequality we used the fact that $d_1 \leq (w/c)^{1/3}$.

Using the fact that $1-x \geq e^{-x-x^2}$ for all $x \in [0, 1/2]$, we can
write
\begin{equation*}
F_1 \geq \prod_{j=1, j\neq i}^{d_1} e^{-x_{i,j} - x_{i, j}^2}.
\end{equation*}
Hence,
\begin{align*}
F_1 \cdot F_2 \geq \exp\left(i^2/(2d_1) - \sum_{j=1, j\neq i}^{d_1} -x_{i,j}
- x_{i, j}^2\right).
\end{align*}

From Equations \eqref{eq:ieq} and \eqref{xijbound}, we know that
\begin{align*}
\sum_{j=1, j\neq i}^{d_1} x_{i,j} + x_{i, j}^2 \leq \frac{c i^2}{w} \cdot
\left(3d_1 + 3 \sqrt{2} i \ln(i+1)\right) \leq \frac{ci^2}{w} \cdot \left(4
d_1 \ln(d_1)\right).
\end{align*}
Hence,
\begin{align*}
F_1 \cdot F_2 & \geq \exp\left(i^2/(2d_1) - \frac{ci^2}{w} \cdot 4 c
\ln(d_1)\right) \\
&= \exp\left(i^2 \left(\frac{1}{2d_1} - \frac{4c^2 \ln(d_1)}{w}\right)\right)
\\
& \geq \exp\left(i^2 \cdot \frac{1}{2d_1} \cdot (1-o(1))\right) \\
& \geq 1.
\end{align*}

\Cref{provedintheappendix} follows.



%

\begin{thebibliography}{57}
\bibitem[8]{bbcmw} Robert Beals, Harry Buhrman, Richard Cleve, Michele Mosca, and Ronald de Wolf. Quantum lower bounds by polynomials. Journal of the ACM, 48(4):778–797, 2001. Earlier version in FOCS’1998, pp. 352-361. quant-ph/9802049. doi:10.1145/502090.502097.
\bibitem[13]{BHMT02} Gilles Brassard, Peter Høyer, Michele Mosca, and Alain Tapp. Quantum amplitude amplification and estimation, volume 305, pages 53–74. American Mathematical Society, 2002. doi: 10.1090/conm/305/05215.
\bibitem[15]{BHT98} Gilles Brassard, Peter Høyer, and Alain Tapp. Quantum cryptanalysis of hash and claw- free functions. In LATIN’98: Theoretical Informatics, pages 163–169, 1998. doi:10.1007/ BFb0054319.
\bibitem[17]{bkt} Mark Bun, Robin Kothari, and Justin Thaler. The polynomial method strikes back: tight quantum query bounds via dual polynomials. In Proceedings of the 50th Annual ACM SIGACT Symposium on Theory of Computing, STOC 2018, Los Angeles, CA, USA, June 25-29, 2018, pages 297–310, 2018. doi:10.1145/3188745.3188784.
\bibitem[18]{bt13} Mark Bun and Justin Thaler. Dual lower bounds for approximate degree and Markov- Bernstein inequalities. In Automata, Languages, and Programming, pages 303–314, 2013. doi:10.1007/978-3-642-39206-1\_26.
\bibitem[20]{coppersmith-rivlin} Don Coppersmith and T. J. Rivlin. The growth of polynomials bounded at equally spaced points. SIAM J. Math. Anal., 23(4):970–983, July 1992. doi:10.1137/0523054.
\bibitem[22]{ez} H. Ehlich and K. Zeller. Schwankung von Polynomen zwischen Gitterpunkten. Mathematische Zeitschrift, 86:41–44, 1964. doi:10.1007/BF01111276.
\bibitem[35]{markov1890question} Andrei Andreyevich Markov. On a question by DI Mendeleev. Zapiski Imperatorskoi Akademii Nauk, 62:1–24, 1890.
\bibitem[39]{mp} Marvin Minsky and Seymour A. Papert. Perceptrons (2nd edition). MIT Press, 1988. First appeared in 1968.
\bibitem[42]{odonnell} Ryan O’Donnell. Analysis of Boolean Functions. Cambridge University Press, USA, 2014.
\bibitem[43]{paturi} Ramamohan Paturi. On the degree of polynomials that approximate symmetric Boolean functions. In Proceedings of the twenty-fourth annual ACM symposium on Theory of computing - STOC 1992, pages 468–474, 1992. doi:10.1145/129712.129758.
\bibitem[46]{rc} T. J. Rivlin and E. W. Cheney. A comparison of uniform approximations on an interval and a finite subset thereof. SIAM Journal on Numerical Analysis, 3(2):311–320, 1966. doi: 10.1137/0703024.
\bibitem[47]{STT12} Rocco Servedio, Li-Yang Tan, and Justin Thaler. Attribute-efficient learning andweight-degree tradeoffs for polynomial threshold functions. In Proceedings of the 25th Annual Conference on Learning Theory, volume 23 of Proceedings of Machine Learning Research, pages 14.1–14.19, 2012. URL: http://proceedings.mlr.press/v23/servedio12.html.
\bibitem[51]{shi} Yaoyun Shi. Approximating linear restrictions of boolean functions. Available at http: //web.eecs.umich.edu/~shiyy/mypapers/, 2002.
\bibitem[55]{Tha16} Justin Thaler. Lower Bounds for the Approximate Degree of Block-Composed Functions. In 43rd International Colloquium on Automata, Languages, and Programming (ICALP 2016), volume 55 of Leibniz International Proceedings in Informatics (LIPIcs), pages 17:1–17:15. Schloss Dagstuhl–Leibniz-Zentrum fuer Informatik, 2016. doi:10.4230/LIPIcs.ICALP.2016. 17.
\bibitem[57]{zhandry} Mark Zhandry. How to construct quantum random functions. In Proceedings of the 2012 IEEE 53rd Annual Symposium on Foundations of Computer Science, FOCS ’12, pages 679–687. IEEE, 2012. doi:10.1109/FOCS.2012.37.
\end{thebibliography}
\end{document}
